\documentclass[11pt,a4paper,openany]{book}
\usepackage[T2A]{fontenc}
\usepackage[utf8]{inputenc}
\usepackage[russian]{babel}
\usepackage{amsmath,amssymb,amsthm,mathtools}
\usepackage{geometry}
\geometry{left=2.4cm,right=3.2cm,top=2.4cm,bottom=2.4cm,marginparwidth=1.9cm,marginparsep=4mm}
\usepackage{enumitem}
\usepackage{tikz}
\usetikzlibrary{patterns,calc,arrows.meta,positioning}
\usepackage[most]{tcolorbox}
\usepackage{xcolor}
\usepackage{marginnote}
\usepackage{booktabs}
\usepackage[hidelinks]{hyperref}

% --- обозначения ---
\newcommand{\ov}[1]{\overline{#1}}
\renewcommand{\emptyset}{\varnothing}
\newcommand{\R}{\mathbb R}
\newcommand{\N}{\mathbb N}
\newcommand{\Z}{\mathbb Z}
\newcommand{\Q}{\mathbb Q}
\newcommand{\M}{\mathfrak M}
\newcommand{\maj}{\operatorname{maj}}
\newcommand{\sd}{\mathbin{\triangle}}
\newcommand{\cupl}{\mathbin{\cup}}
% номер страницы оригинала на полях
\newcommand{\bookpage}[1]{\leavevmode\marginnote{\footnotesize\textcolor{gray}{[с.~#1]}}\ignorespaces}
% подпись к рисункам в стиле книги
\newcommand{\figcap}[1]{\par\smallskip\centering\small #1\par}
% --- добавленный материал ---
\definecolor{solcol}{RGB}{20,80,150}
\newtcolorbox{solution}[1][]{breakable,enhanced,colback=solcol!4,colframe=solcol!60,
  boxrule=0.5pt,arc=1pt,left=6pt,right=6pt,top=4pt,bottom=4pt,
  title={Решение упражнения #1},fonttitle=\bfseries\small,coltitle=white,colbacktitle=solcol!75}
\newtcolorbox{remarkbox}[1][]{breakable,enhanced,colback=orange!5,colframe=orange!60!black,
  boxrule=0.4pt,arc=1pt,title={#1},fonttitle=\bfseries\small}
% исправление опечатки: отметка в тексте
\newcommand{\fixed}[1]{\footnote{\textit{Исправлено:} #1}}
\theoremstyle{plain}
\newtheorem*{lemma*}{Лемма}
\newtheorem*{theorem*}{Теорема}
\newtheorem*{prop*}{Утверждение}
\newcommand{\exitem}[1]{\par\smallskip\noindent\textbf{#1}~}

\emergencystretch=2.5em
\usepackage{fancyhdr}
\pagestyle{fancy}\fancyhf{}
\fancyhead[LE,RO]{\thepage}\fancyhead[RE]{\small\itshape В.\,Л.~Рвачев. Теория $R$-функций}\fancyhead[LO]{\small\itshape\nouppercase{\rightmark}}
\fancypagestyle{plain}{\fancyhf{}\fancyfoot[C]{\thepage}\renewcommand{\headrulewidth}{0pt}}

\title{В.\,Л.~Рвачев\\[4pt]
\Large Теория $R$-функций и некоторые ее приложения\\[2pt]
\large (Киев: Наукова думка, 1982)\\[14pt]
\normalsize Страницы 1--29: предисловие, введение, гл.~1, \S\S\,1--4\\
текст с исправленными опечатками и решениями упражнений}
\author{}
\date{}

\begin{document}
\maketitle
\thispagestyle{empty}
\section*{О настоящем издании}
Это набор в \LaTeX\ первых 29 страниц монографии В.\,Л.~Рвачева
(предисловие, введение и \S\S\,1--4 гл.~1), выполненный по скану
книги (файл \texttt{r-functions-theory-rvachev.pdf} в репозитории).

\begin{itemize}[nosep]
\item Номера страниц оригинала указаны на полях серым цветом: \textcolor{gray}{[с.~15]}.
\item Нумерация формул, рисунков, упражнений сохранена.
\item Замеченные опечатки оригинала исправлены прямо в тексте; каждое
  исправление отмечено сноской, а полный список приведен в приложении.
\item Решения упражнений, отсутствующие в книге, вставлены сразу после
  соответствующего блока <<Упражнения и задачи>> в рамках
  \begin{tcolorbox}[colback=solcol!4,colframe=solcol!60,boxrule=0.5pt,arc=1pt,
    left=3pt,right=3pt,top=1pt,bottom=1pt,width=5.3cm,nobeforeafter,box align=center]
  \small такого вида\end{tcolorbox}.
  Собственные решения автора (там, где они есть) оставлены как часть основного текста.
\item Рисунки перерисованы (TikZ) по оригиналу.
\end{itemize}
\clearpage

\tableofcontents

\chapter*{Предисловие}
\addcontentsline{toc}{chapter}{Предисловие}\markboth{Предисловие}{Предисловие}
\bookpage{3}

Повышенный интерес к теории $R$-функций (или, если говорить с более общих
позиций, теории $R$-отображений) объясняется как простотой и привлекательностью
ее начальных идей, так и широким диапазоном областей, в которых она находит
приложение. В этом отношении теория $R$-функций напоминает теорию групп: обе
эти теории базируются на чисто алгебраических концепциях, а их приложения
весьма многообразны. Для $R$-функций областями применения оказались:
аналитическая геометрия, математическое программирование, оптимальное
размещение геометрических объектов, распознавание образов, теория
устойчивости, химическая технология и в особенности математическая физика
с ее многочисленными задачами исследования, расчета и оптимизации
физико-механических полей (температурных, деформационных, силовых,
электромагнитных, магнитогидродинамических и~др.).

Цель монографии "--- изложить основные результаты, полученные почти за
двадцатилетний период в теории $R$-функций и ряде ее приложений. Особенно
интенсивно стала развиваться и применяться теория $R$-функций в последние
несколько лет. В связи с этим возникла необходимость осмысливания и обобщения
полученных результатов, однако вышедшие в свет монографии в основном были
посвящены отдельным прикладным направлениям и не отражали в должной мере
прогресс в развитии многих вопросов теории.

Монография состоит из введения и четырех глав. Гл.~1 вводит читателя в круг
представлений алгебры логики и исчисления предикатов (называемых также
дискретной или конечной математикой). Эти вопросы освещаются лишь в объеме,
необходимом для дальнейшего изложения. Более полно данные вопросы рассмотрены
в работе С.\,В.~Яблонского~[194]. В гл.~2 вводятся понятия $R$-отображения
($R$-функции), чертежа, уравнения чертежа, алгебраически полной системы и так
далее (термин <<чертеж>> используется вместо часто употребляемых терминов
<<линия>>, <<фигура>>, <<поверхность>>, <<тело>>, <<гиперповерхность>>,
<<геометрический объект>> и~т.\,п.). В этой же главе решается принципиально
важный вопрос о построении уравнений чертежей (решение обратной задачи
аналитической геометрии), рассматриваются приложения полученных результатов
в математическом программировании и оптимальном размещении геометрических
объектов. Новый в идейном плане материал в основном содержится в гл.~3 и~4.

Материал гл.~3, в которой используются результаты предыдущих глав, посвящен
развитию конструктивных средств математики и решению казавшейся ранее
неразрешимой проблемы построения полных координатных последовательностей,
удовлетворяющих различного типа краевым условиям для областей
\bookpage{4}%
практически произвольной формы. Эти результаты легли в основу новой
технологии программирования и ее реализации в виде программирующих систем
(генераторов программ) серии <<Поле>>, позволивших в десятки раз сократить
затраты на составление и отладку программ. О генераторах программ и решении
краевых задач в условиях их эксплуатации речь идет в гл.~4. Материал этой
главы подобран с таким расчетом, чтобы не только показать возможности
применения $R$-функций к решению задач расчета полей различного типа, но и
отметить некоторые специфические особенности реализации теории $R$-функций на
практике.

Многие параграфы сопровождаются упражнениями и задачами. Некоторые из них
носят проблемный характер, и автор не знает их решений. Ряд нерешенных
вопросов отмечен в основном тексте. Большое число примеров, приведенных
в работе, облегчает освоение материала, делает его доступным для студентов
старших курсов инженерных и математических специальностей. (Примеры
и некоторые параграфы, которые при первом чтении можно опустить, набраны
петитом.) В списке литературы указаны лишь те источники, содержание которых
в какой-то мере связано с излагаемыми результатами. Обширная библиография
содержится в монографиях и статьях автора и его учеников И.\,В.~Гончарюка,
В.\,С.~Проценко, Ю.\,Г.~Стояна, В.\,А.~Рвачева, А.\,П.~Слесаренко,
Г.\,П.~Манько, Т.\,И.~Шейко, Н.\,С.~Синекопа, Ф.\,Ф.~Коваля и~др. [25--27, 49,
50, 67--75, 100, 101, 109--148, 154--156, 160--164, 167--172, 184--188].
Ссылки (на теоремы, примеры, упражнения и~т.\,д.) внутри параграфов (глав)
даются без указаний номеров параграфов (глав).

\chapter*{Введение}
\addcontentsline{toc}{chapter}{Введение}\markboth{Введение}{Введение}
\bookpage{5}

Существует обширный круг научно-технических проблем, при решении которых
необходимо учитывать и преобразовывать геометрическую информацию.
В простейших случаях (например, при задании информации о машиностроительном
чертеже для хранения ее в памяти ЭВМ) вполне можно ограничиться тем или иным
способом кодирования элементов чертежа или в крайнем случае поточечным их
сканированием. Такие задачи в большинстве случаев носят чисто технический
характер и их решение не сопряжено с преодолением принципиальных затруднений.
Вместе с тем во многих случаях простого кодирования геометрической
информации недостаточно, а необходимо описание ее в той форме, какая принята
в аналитической геометрии, т.\,е. с помощью уравнений или неравенств. Речь
идет о том, чтобы так же, как, например, для прямой на плоскости или шара
в пространстве, можно было бы и для таких <<сложных>> геометрических
объектов, как усеченный конус с цилиндрическим отверстием, ломаная линия,
зубчатое колесо и других применять уравнения $\omega(x)=0$, $x\in\mathbb R^n$
(или неравенства $\omega(x)\ge0$). Здесь $\omega(x)$ "--- функция, которая
имеет вид единого аналитического выражения и обладает, возможно, еще
и некоторыми дополнительными свойствами.

Естественно, что при такой общей постановке возникает ряд вопросов,
требующих строгой формализации (и такая формализация выполнена в данной
работе). К их числу относятся следующие: что такое сложный геометрический
объект и единое аналитическое выражение, при каких условиях поставленная
задача разрешима, имеются ли эффективные алгоритмы решения этой задачи
и~т.\,д.?

В историческом плане сформулированная выше проблема восходит еще к Декарту
и известна в литературе как обратная задача аналитической геометрии: задан
геометрический объект, требуется написать его уравнение. (Прямая задача:
дано уравнение $\omega(x)=0$, $x\in\mathbb R^n$, найти соответствующий ему
в~$\mathbb R^n$ геометрический объект $L$, т.\,е. совокупность всех точек
$\mathbb R^n$, которые удовлетворяют уравнению $\omega(x)=0$.) Во времена
Декарта запас изученных линий и поверхностей (чертежей) мало чем отличался от
того, каким располагали математики древности. Метод координат
\bookpage{6}%
позволил беспредельно увеличивать число изученных чертежей, так как каждое
новое уравнение давало и новый чертеж. Это привело к тому, что в дальнейшем
основное внимание было уделено изучению чертежей, описываемых заданными
уравнениями. Обратная же задача рассматривалась лишь для простейших
геометрических объектов: прямой, конических сечений, эллипсоидов,
параболоидов и~др.

В аналитической геометрии обычно ограничиваются системой
$H_0=\{x_1+x_2,\ x_1x_2,\ a=\mathrm{const}\ \forall a\in\mathbb R\}$
операций, применяемых для написания уравнений, которая приводит к множеству
$\mathfrak M(H_0)$ целых рациональных функций (полиномов). Ему соответствует
множество $\mathfrak R(H_0)$ чертежей, описываемых уравнениями $\omega(x)=0$,
$\omega(x)\in\mathfrak M(H_0)$ и называемых алгебраическими. Однако множество
алгебраических чертежей все же весьма <<бедно>>. Оно не содержит, например,
таких чертежей, как прямоугольник, двутавр, поверхность искусственного
спутника Земли и~т.\,п. Правда, существует принципиальная возможность
аппроксимации сложных чертежей алгебраическими (это следует из известной
теоремы Вейерштрасса о приближении непрерывных функций полиномами~[7]), но
сама по себе реализация ее весьма сложная проблема. То же можно сказать
и относительно известной теоремы о возможности сколь угодно хорошей
аппроксимации замкнутых кривых на плоскости лемнискатами (т.\,е. также
алгебраическими кривыми)~[150].

Существенным расширением множества $\mathfrak R(H_0)$ алгебраических
чертежей является множество $\mathfrak R_{1/2}(H_0)$ полуалгебраических
чертежей~[5]. Оно включает в себя всякие чертежи, которые также могут быть
описаны с помощью упомянутой системы операций $H_0$, но уже не одним
уравнением, а системами систем уравнений и неравенств вида $f_{ik}=0$,
$\varphi_{jk}\ge0$, где $f_{ik},\varphi_{jk}\in\mathfrak M(H_0)$, $i$, $j$,
$k$ "--- принадлежат конечным множествам индексов (точка принадлежит
чертежу, если хотя бы одна из систем удовлетворяется). Из определения
множества $\mathfrak R_{1/2}(H_0)$ следует, что
$\mathfrak R(H_0)\subset\mathfrak R_{1/2}(H_0)$. Нетрудно установить также
(см. \S\,9 гл.~2), что существуют полуалгебраические, но не алгебраические
чертежи, т.\,е. $\mathfrak R(H_0)\ne\mathfrak R_{1/2}(H_0)$.

Пусть $H$ есть некоторое расширение системы операций $H_0$. Тогда, подобно
тому как системе $H_0$ были поставлены в соответствие множества чертежей
$\mathfrak R(H_0)$ и $\mathfrak R_{1/2}(H_0)$, системе $H$ будут
соответствовать множества $\mathfrak R(H)$ и $\mathfrak R_{1/2}(H)$
чертежей, которые назовем $H$-реализуемыми и полу-$H$-реализуемыми
соответственно. Правомерен вопрос: нельзя ли так выбрать систему $H$, чтобы
получить равенство $\mathfrak R(H)=\mathfrak R_{1/2}(H)$? Располагая такой
системой, можно было бы задавать сложные геометрические объекты не
упомянуты\-%
\bookpage{7}%
ми выше системами систем уравнений и неравенств, а одним уравнением,
написанным в символах системы $H$. В гл.~2 будет показано, что такие
системы, названные алгоритмически полными, существуют и могут быть
построены с помощью $R$-функций. Впервые понятие $R$-функции как раз и было
введено в связи с возникшей необходимостью составлять уравнения сложных
чертежей (в процессе поиска приближенных решений некоторых пространственных
контактных задач теории упругости).

Со временем понятие $R$-функции было обобщено и возникло общее понятие
$R$-отображения (см. \S\,2 гл.~2). В основе этого понятия лежит выделение из
систем алгебраических операций (т.\,е. отображений вида
$f\colon\mathfrak X^n\to\mathfrak X$, где $\mathfrak X$ "--- некоторое
множество) таких, которые в некотором смысле соответствуют функциям
$k$-значной логики (т.\,е. алгебраическим операциям вида
$F\colon B_k^n\to B_k$, у которых множество $B_k$ состоит из $k$
элементов). В общем виде теория $R$-отображений еще мало разработана.
В основном изучались $R$-функции, которые строятся на множестве
$\mathfrak X=\mathbb R=(-\infty,\infty)$, и соответствуют функциям
двузначной (булевым функциям) и трехзначной логики. Именно эти функции
позволили решить обратную задачу аналитической геометрии, а впоследствии
найти и другие приложения на более высоких уровнях прикладной математики.
При этом оказалось, что существуют элементарные алгоритмически полные
системы и в целом само множество основных элементарных функций (в сочетании
с арифметическими операциями) алгоритмически полное. Это означает,
в частности, что всякий чертеж, составленный из <<кусков>> элементарных
линий или поверхностей, сам элементарен, т.\,е. может быть описан
уравнением $\omega(x)=0$, где $\omega(x)$ "--- элементарная (имеет вид
единого аналитического выражения) функция.

В первых работах по теории $R$-функций был определен и достаточно широкий
круг возможных ее приложений: оптимальное размещение геометрических объектов
(в частности, оптимальный раскрой [109, 142, 143]), распознавание образов
[157, 158], математическое программирование [109], конструктивная теория
функций [61, 111] и в особенности краевые задачи математической физики [111,
128, 134].

Одно из неожиданных направлений, в котором $R$-функции нашли применение,
"--- это теория устойчивости движения. Первоначально метод $R$-функций
применен в этом направлении при исследовании области устойчивости
неустановившегося движения~[77]. За истекшее время с использованием
$R$-функций получены новые результаты при численной реализации классических
критериев устойчивости [78, 79, 82], построении функций Ляпунова, синтезе
систем, асимптотически устойчивых по первому приближению~[79], выборе
параметров функций связей крупномасштабных систем с лишь устойчивыми
подсистемами~[80], в задаче об устойчивости аналитических движений~[76]
и других~[77]. Этим, судя по всему, не исчерпываются возможности применения
$R$-функций. В част\-%
\bookpage{8}%
ности, представляют значительный интерес разработка теории
$R$-отображений в пространстве параметров динамических систем~[81],
построение и исследование свойств $R$-отображений, управляемых специальными
системами сравнения [76--82], а также применение $R$-функций в задачах
распознавания образов [157, 158], химической технологии~[46], медицинской
диагностики~[19] и других направлениях [7, 36].

Математическому программированию и размещению геометрических объектов
посвящены \S\S\,26, 27 гл.~2. В основном они написаны на уровне
монографии~[109], изданной в 1967~г., более глубоко эти вопросы освещены
в работах Ю.\,Г.~Стояна и его учеников [167--172]. Кратко изложены также
вопросы, относящиеся к теории сплайнов и атомарных функций (\S\S\,15, 16
гл.~3). Более детально эти вопросы рассмотрены в монографии~[127]. Основные
результаты в развитии теории атомарных функций принадлежат В.\,А.~Рвачеву
[123--127], первые приложения "--- Т.\,И.~Шейко~[184].

В последние годы в индустриально развитых странах возросли темпы создания
новых математических теорий и разработки на их основе методов решения на ЭВМ
задач исследования, расчета и оптимизации физико-механических полей. Все
возрастающий интерес к проблеме физико-механических полей объясняется тем,
что ею охвачен широкий круг направлений, таких как электродинамика,
теплофизика, теория фильтрации, теория упругости и пластичности, магнитная
гидродинамика и другие, развитие которых имеет первостепенное значение для
научно-технического прогресса. Математическими моделями физико-механических
полей являются задачи для уравнений с частными производными при
определенных краевых и начальных условиях. Специфическая особенность полей
"--- их зависимость не только от характера физических законов, учитываемых
соответствующими уравнениями, но и от формы, взаимного расположения тел,
в которых поля возбуждаются, конфигурации площадок их контактного
взаимодействия и других геометрических и физических факторов.

Наличие в постановке краевых задач двух разнородных видов информации "---
аналитической и геометрической "--- серьезное препятствие при создании
методов и алгоритмов решения этих задач. Дело в том, что всякий метод
неизбежно должен учитывать оба упомянутые вида информации, а это требует
преобразования геометрической информации в аналитическую (при этом не всякое
кодирование геометрической информации годится). В таких классических
методах, как, например, разделения переменных и интегральных преобразований,
форма областей и участков их границ учитывается благодаря подходящему выбору
систем координат, в методе конформных отображений "--- при построении
отображающих функций, в вариационных методах "--- при построении
координатных функций, в сеточных методах "--- при составлении уравнений
для узлов, близких к границе, и~т.\,д. Метод конечных элементов
(МКЭ)~[31, 66] возник именно в связи со стремлением как
\bookpage{9}%
можно точнее учитывать геометрические формы. Появился также ряд методов
(Рейснера, штрафов и~др.), в которых краевые условия естественны и
формально им можно не удовлетворять [58, 88]. В этих методах геометрическая
информация учитывается путем интегрирования по областям соответствующей
формы (т.\,е. геометрическая информация учитывается интегрально).

Особенно актуальна разработка таких методов решения краевых задач, которые
имели бы универсальный характер и не требовали от исследователя (как
правило, не математика) знания тонких вопросов теории. Кроме того,
универсальность обусловливает возможность привлечения методов системного
программирования, что имеет существенное значение для автоматизации научных
исследований в области краевых задач. Конечно, для того или иного
конкретного класса задач существуют или могут быть созданы специальные
методы, превосходящие по эффективности любой универсальный метод, но их
использование требует, как правило, хорошей математической и специальной
подготовки и значительных затрат времени, что может оказаться неприемлемым
для быстрого решения все новых и новых задач, различных по специфике. Из
упомянутых выше методов универсальны разностные, вариационные и конечных
элементов. Можно назвать также методы Трефтца, локальных вариаций~[183],
штрафов~[112] и~др. Перечисленные методы принято называть прямыми; для них
характерно сведение краевых задач к решению алгебраических систем линейных
уравнений. Есть, однако, другая, более глубокая связь между этими методами,
которая делает различие между ними весьма условным: во всех прямых методах
приближенное решение краевой задачи отыскивается в виде
\begin{equation}\tag{1}\label{eq:intro1}
u_N=\sum_{k=1}^{N}C_k\varphi_k(x)+\varphi_0(x),
\end{equation}
где $\varphi_k(x)$ "--- некоторые заранее выбранные функции, а $C_k$ "---
числа, которые находятся в результате решения упоминавшихся выше систем
алгебраических уравнений. Заметим, что в некоторых методах (сеточных, МКЭ)
формула~(1) представления решения не всегда явно выписывается, однако это
не меняет сути. В простейших конечно-разностных методах, например, за основу
берется функция $y=h_1(t)$ (финитный сплайн Шенберга первого порядка).

Приближенное решение (например, для двумерной задачи) выбирается в виде
линейной комбинации произведений сдвигов сжатий функций~$h_1$:
\begin{equation}\tag{2}\label{eq:intro2}
u_N(x_1,x_2)=\sum_{i,j}C_{ij}\,h_1(\alpha x_1-\beta i)\,h_1(\gamma x_2-\delta j).
\end{equation}
Здесь $C_{ij}=u_{ij}$ "--- значение искомой функции в узле $(i,j)$, а
$\alpha$, $\beta$, $\gamma$ и $\delta$ выбираются такими, чтобы функция
$\varphi_{ij}=h_1(\alpha x_1-\beta i)\,h_1(\gamma x_2-\delta j)$ была равна
единице в $(i,j)$-м узле и
\bookpage{10}%
нулю "--- в остальных узлах сетки. Нетрудно представить, что графически
функция~(2) "--- это многогранная поверхность. При этом изменение
какой-либо из величин $C_{ij}$ приводит лишь к локальному ее изменению
в окрестности узла $(i,j)$. Если, учтя граничные условия, подставить
формулу~(2) в функционал, соответствующий краевой задаче, и написать условия
его экстремума, то относительно величин $C_{ij}$ получим систему
конечно-разностных уравнений. Другие конечно-разностные схемы получаем
аналогично, используя иные финитные функции (например, сплайны Шенберга
второго и третьего порядков) (см. \S\,15 гл.~3), что соответственно
приводит к менее разреженным матрицам систем алгебраических уравнений.
В этом смысле ничем по существу не отличается от сеточных методов и метод
конечных элементов, в котором функции $\varphi_k(x)$ также локальные сплайны
(в большинстве случаев дефектные, т.\,е. порядок используемых в них
полиномов превосходит их гладкость больше чем на единицу), однако их
носители могут иметь более сложную геометрическую форму, разную в различных
подобластях рассматриваемой области. Именно вследствие этого
и осуществляется более гибкая привязка к геометрии области.

В классических вариационных методах функции $\varphi_k(x)$ "--- это так
называемые координатные функции, вообще, не финитные, удовлетворяющие, если
это необходимо, краевым условиям. Применение таких функций приводит
к системе алгебраических уравнений, как правило, с полностью заполненными
матрицами. Такие системы обычно плохо обусловлены даже при сравнительно
небольших $N$ (порядка $10^2$) и требуют для хранения элементов матрицы
в памяти ЭВМ $N^2$ ячеек. Кроме того, вычисление коэффициентов в методах
типа Ритца "--- Галеркина или наименьших квадратов требует вычисления
кратных интегралов по всей области, в которой решается задача, и как
следствие "--- больших затрат машинного времени. Таким образом, если взять
за основу один и тот же вариационный принцип (например, минимизировать
функционал Ритца), то метод сеток, конечных элементов и классический
вариационный будут отличаться лишь выбором функций $\varphi_k(x)$. Будут
различными при одном и том же $N$ вычислительные затраты и накопление
погрешностей, что, конечно, очень важно, однако тип сходимости, вообще,
будет один и тот же.

Возможен и другой вариант: функции $\varphi_k(x)$ выбираются одни и те же,
но по-разному определяются постоянные $C_k$. Например, используя одну и ту
же формулу~(1), можно находить постоянные по методу Ритца, а можно, добавив
к функционалу Ритца соответствующие слагаемые, применить метод Куранта
и получить благодаря этому более высокий тип сходимости, чем сходимость по
энергии.

Таким образом, все прямые методы объединяет между собой формула~(1),
а отличие между ними состоит лишь в выборе функций $\varphi_k(x)$
и способах отыскания постоянных $C_k$.
\bookpage{11}%

Каким же условиям должны удовлетворять функции $\varphi_k(x)$, если
проблему рассматривать с позиций теории приближений, вычислительной
математики и системного программирования? В большинстве случаев требуется,
а в остальных желательно, чтобы функции $\varphi_k(x)$ удовлетворяли краевым
условиям рассматриваемой задачи. Основная трудность в решении этого вопроса
связана с тем, что в формулы, представляющие функции $\varphi_k(x)$, должна
некоторым образом включаться информация геометрического характера о форме
границы рассматриваемой области (и форме ее участков для краевых условий
смешанного типа). Как уже отмечалось, в некоторых методах (сеточных, МКЭ)
эта информация учитывается приближенно вследствие локального приближения
узлов сетки к границе, а в некоторых (Рейснера, штрафов) "--- в интегральном
смысле. Это, конечно, накладывает определенный отпечаток на тип сходимости
и величину погрешностей. Тем не менее с этими недостатками можно было бы
мириться при условии, что погрешность в результате приближенного учета
геометрических форм соразмерна с погрешностью самого приближенного метода.
Однако есть еще один весьма важный фактор, который необходимо учитывать:
и в сеточных методах, и в МКЭ составление уравнений для пограничных узлов
трудно формализуемое, что затрудняет решение вопросов автоматизации
программирования.

Рассмотрим, какая ситуация складывается в методе штрафов. В простейшем
варианте этого метода краевая задача (3.263), (3.264) (см. \S\,9 гл.~3)
заменяется задачей о минимуме функционала
\begin{equation}\tag{3}
I=\int_\Omega(Au-f)^2\,d\Omega+\sum_{i=1}^{m}M_i\int_{\partial\Omega_i}(L_iu-\varphi_i)^2\,d\Gamma_i,
\end{equation}
где $M_i>0$ "--- некоторые постоянные. Нетрудно заметить, что минимум $I$
достигается на решении (обобщенном) этой краевой задачи. Однако при
численной реализации метода штрафов не ясно, как следует выбирать
постоянные $M_i$. Ответ можно получить на основе численных экспериментов,
но организация этих экспериментов "--- дело сложное, требующее больших
затрат времени. (Видимо, этим и объясняется то обстоятельство, что
к настоящему времени опубликованы лишь немногие решения реальных краевых
задач, полученные методом штрафов, и касаются они лишь простых случаев.)
Поэтому, если бы появилась возможность строить в конструктивном отношении
достаточно простые функции $\varphi_k(x)$ так, чтобы краевые условия
удовлетворялись точно и в то же время были сохранены все другие
положительные факторы, характерные для упомянутых методов, то это имело бы
важное значение для их дальнейшего развития.

В течение длительного времени задача построения в явном виде функций
$\varphi_k(x)$, точно удовлетворяющих произвольным краевым условиям для
областей сложной формы, казалась практически неразрешимой. Дж.~Ортега
и В.~Рейнболдт [95, с.~36] по этому поводу,
\bookpage{12}%
например, пишут: <<Хотя формально метод Ритца распространяется на
двумерные вариационные задачи и задачи более высокой размерности, здесь
имеется серьезная практическая трудность, состоящая в построении подходящих
базисных функций для общих областей>>.

В гл.~3, 4 показано, что эта проблема может быть решена с помощью
конструктивного аппарата теории $R$-функций. При этом речь идет не просто об
однородной задаче Дирихле, а о краевых условиях самых различных типов для
областей практически произвольной формы. Для линейных краевых условий
получаются формулы вида~(1), а для нелинейных "--- формулы общего вида:
\begin{equation}\tag{4}
u_N=f(x,C_1,\dots,C_N).
\end{equation}

По теории $R$-функций "--- довольно много работ (количество их с каждым
годом возрастает), в том числе монографии [111, 127, 128, 134, 135].
Заметим, что в этой теории речь идет не только об удовлетворении краевым
условиям. С точки зрения теории приближений одного лишь удовлетворения
краевым условиям не достаточно для того, чтобы можно было получать
приближения к точному решению. Задачу можно считать успешно решенной лишь
при обеспечении определенных аппроксимационных свойств функций
$\varphi_k(x)$: их линейные комбинации~(1) должны достаточно хорошо
приближать функции того класса, которому принадлежит точное решение задачи,
т.\,е. функции~(1) должны удовлетворять условиям полноты, а по возможности
и условию аппроксимационной универсальности~[126].

Основное внимание в данной монографии обращено на конструктивную сторону
теории $R$-функций. Именно эта сторона сыграла главную роль при создании
генераторов программ серии <<Поле>> (\S\S\,1, 2 гл.~4), позволивших во много
раз сократить затраты на составление и отладку программ при решении краевых
задач. (Более подробная информация содержится в публикациях Г.\,П.~Манько
[67--71].)

Развитие конструктивного начала в математике всегда происходило наряду
с развитием самой теории, причем вызывалось оно в первую очередь запросами
практики. Например, если бы в свое время не был введен символ $\tg x$,
а каждый раз вместо него писали бы $\dfrac{\sin x}{\cos x}$, то от этого не
пострадала бы ни одна из теорем математики. Но нетрудно представить,
насколько усложнилось бы решение многих прикладных проблем.

Функция $\tg x$ выдержала испытание <<естественным отбором>>. В связи
с появлением ЭВМ темпы этого <<отбора>> неизмеримо возросли. Одна из задач
данной работы "--- показать, что есть новые функции (элементарные
и неэлементарные), которые по широте применения на практике и важности для
развития теории, несомненно, выдержат испытание временем.


\chapter{Элементы дискретной математики}
\bookpage{13}
\section{Множества и логические операции над ними}

Логические операции над множествами в настоящее время применяются не только
в математике, но и в многочисленных ее приложениях. В данной работе им
отводится особая роль: формализация, достигаемая на уровне алгебры
множеств, будет представлять собой существенное звено в конструктивных
построениях, осуществляемых с помощью $R$-функций (в частности, при
построении уравнений сложных геометрических объектов и структур решений
краевых задач). Заметим, что в ряде случаев наряду с термином
<<множество>> используются его синонимы <<совокупность>>,
<<многообразие>>, <<система>> и~др.

Примем следующие обозначения:

\begin{itemize}[nosep,label={},leftmargin=1em]
\item $\N$ "--- множество неотрицательных целых чисел;
\item $\Z$ "--- множество произвольных целых чисел;
\item $\Q$ "--- множество рациональных чисел;
\item $\R$ "--- множество действительных чисел;
\item $C^m(\Omega)$ "--- множество $m$ раз непрерывно дифференцируемых функций;
\item $\emptyset$ "--- пустое множество;
\item $2^A$ "--- множество всех подмножеств множества $A$;
\item $f(A)$ "--- образ множества $A$ при отображении $f$;
\item $f^{-1}(A)$ "--- прообраз множества $A$ при отображении $f$;\fixed{в оригинале <<при отображении $|$>>.}
\item $f/A$ "--- сужение (ограничение) отображения $f$ на множество $A$
  (т.\,е. $f/A$ есть отображение, совпадающее с $f$ на множестве $A$. Вне
  множества $A$ отображение $f/A$, вообще, не определено);
\item $\inf A$ ($\sup A$) "--- нижняя (верхняя) грань множества $A\subset\R$;
\item $a\in A$ (или $A\ni a$) "--- элемент $a$ принадлежит множеству $A$
  ($a$ находится в состоянии принадлежности к $A$);
\item $a\notin A$ (или $A\not\ni a$) "--- элемент $a$ не принадлежит
  множеству $A$ ($a$ находится в состоянии непринадлежности к $A$);
\item $A\subset B$ (или $B\supset A$) "--- множество $A$ является
  подмножеством множества $B$;
\item $A\equiv B$ "--- множества $A$ и $B$ тождественны, т.\,е.
  $A\supset B$ и $A\subset B$;
\item $\exists a\in A\dots$ "--- существует (по крайней мере один) элемент
  $a$ множества $A$, что \dots\ ($\exists$ "--- квантор существования);
\bookpage{14}%
\item $\forall a\in A\dots$ "--- для всякого элемента $a$ множества
  $A$\dots\ ($\forall$ "--- квантор общности);
\item $A_1\times\dots\times A_n$ "--- прямое произведение множеств
  $A_1,\dots,A_n$ есть множество упорядоченных строк $a_1,\dots,a_n$, где
  $a_i\in A_i$. В частности, если $A_1=\dots=A_n$, то
  $A\times\dots\times A=A^n$ называется $n$-й степенью множества $A$.
  Например, $n$-мерное евклидово пространство $\R^n$ есть $n$-я степень
  множества~$\R$.
\end{itemize}

Перейдем к логическим операциям над множествами. Операцию $F$ над
множествами $A,B,\dots$ будем называть логической, если в результате ее
применения получается множество, принадлежность элементов к которому
вполне определяется их состояниями принадлежности или непринадлежности
к множествам $A,B,\dots$ Примеры логических операций следующие:

1. $A\cap B$ (или $AB$) "--- \emph{пересечение} множеств $A$ и $B$ "--- есть
множество, состоящее из всех элементов, принадлежащих как множеству $A$,
так и множеству $B$.

2. $A\cup B$ (или $A+B$) "--- \emph{соединение} (сумма) множеств $A$ и $B$
"--- есть множество, состоящее из всех элементов, принадлежащих хотя бы
одному из множеств $A$ или $B$.

3. $A\setminus B$ (или $A-B$) "--- \emph{разность} множеств $A$ и $B$ "---
есть совокупность тех элементов $A$, которые не принадлежат $B$:
\begin{equation}\tag{1.1}
A\setminus B=A\cap\ov B.
\end{equation}

4. $A\sd B$ "--- \emph{симметрическая разность} "--- есть соединение
разностей $A\setminus B$ и $B\setminus A$, т.\,е.
\begin{equation}\tag{1.2}
A\sd B=(A\cap\ov B)\cup(B\cap\ov A).
\end{equation}

Часто приходится рассматривать системы множеств, являющиеся подмножествами
некоторого основного (так называемого унитарного) множества $I$. Для них
вводится понятие дополнения (или отрицания).

5. Множество $\ov A=I\setminus A$ называется дополнением (отрицанием)
множества $A$. (Иногда употребляются обозначения $\ov A=A'={\sim}A$ и~др.)

6. Тождество $F(A,B,\dots)\equiv\Phi(A,B,\dots)$, где $F$ и $\Phi$ "---
некоторые формулы, построенные с помощью тех или иных логических операций
над множествами $A,B,\dots$, означает равенство множеств $F(A,B,\dots)$
и $\Phi(A,B,\dots)$ при любом выборе множеств $A,B,\dots$

Нетрудно убедиться, что операции пересечения, соединения и дополнения
обладают следующими свойствами, выраженными в виде тождеств.

I. Коммутативность:
\[
1^\circ)\ A\cap B\equiv B\cap A;\qquad 2^\circ)\ A\cup B\equiv B\cup A.
\]
II. Ассоциативность:
\[
\begin{aligned}
&3^\circ)\ A\cap(B\cap C)\equiv(A\cap B)\cap C,\\
&4^\circ)\ A\cup(B\cup C)\equiv(A\cup B)\cup C.
\end{aligned}
\]
\bookpage{15}%
III. Дистрибутивность:
\[
\begin{aligned}
&5^\circ)\ A\cap(B\cup C)\equiv(A\cap B)\cup(A\cap C);\\
&6^\circ)\ A\cup(B\cap C)\equiv(A\cup B)\cap(A\cup C).
\end{aligned}
\]
IV. Идемпотентность:
\[
7^\circ)\ A\cap A\equiv A;\qquad 8^\circ)\ A\cup A\equiv A.
\]
V. Закон двойного отрицания:
\[
9^\circ)\ \ov{(\ov A)}\equiv\ov{\ov A}\equiv A.
\]
VI. Правила де Моргана:
\[
10^\circ)\ \ov{A\cap B}\equiv\ov A\cup\ov B;\qquad
11^\circ)\ \ov{A\cup B}\equiv\ov A\cap\ov B.
\]
VII. Закон логического противоречия:
\[
12^\circ)\ A\cap\ov A\equiv\emptyset.
\]
VIII. Закон исключенного третьего:
\[
13^\circ)\ A\cup\ov A\equiv I.
\]
IX. Операции с унитарным множеством $I$ и пустым множеством $\emptyset$:
\[
\begin{array}{ll}
14^\circ)\ I\cap A\equiv A; & 17^\circ)\ \emptyset\cup A\equiv A;\\
15^\circ)\ I\cup A\equiv I; & 18^\circ)\ \ov\emptyset\equiv I;\\
16^\circ)\ \emptyset\cap A\equiv\emptyset; & 19^\circ)\ \ov I\equiv\emptyset.
\end{array}
\]
Приведенные формулы позволяют осуществлять тождественные преобразования
формул, записанных в символах операций $A\cap B$, $A\cup B$, $\ov A$,
например, с целью их упрощения.

{\small\noindent\textbf{Пример.}
\[
(A\cap\ov B\cap C)\cup(A\cap C)\equiv(A\cap C\cap\ov B)\cup(A\cap C\cap I)
\equiv(A\cap C)\cap(\ov B\cup I)\equiv(A\cap C)\cap I\equiv A\cap C.
\]
(Здесь использованы свойства $1^\circ$, $3^\circ$, $5^\circ$, $14^\circ$,
$15^\circ$.)\fixed{в оригинале в левой части стоит
$(A\cap\ov B\cap C)\cap(A\cap C)$; по ходу преобразований (вынесение
$A\cap C$ по закону $5^\circ$) первым знаком должен быть $\cup$.}\par}

Логическая терминология, использованная для наименования некоторых из
приведенных выше формул, имеет свои корни в метаматематике~[84],
в частности в исчислении высказываний или булевой алгебре. О последней
речь будет идти подробнее в \S\S\,6--8 данной главы.

Справедливость свойств $1^\circ$--$19^\circ$ легко проверяется. Убедимся,
например, в справедливости формулы $5^\circ$. Пусть
$x\in A\cap(B\cup C)$, т.\,е. $x$ "--- элемент множества, стоящего в левой
части формулы $5^\circ$. Это означает, что $x$ принадлежит одновременно как
множеству $A$, так и множеству $B\cup C$. Но тогда $x\in B$ или $x\in C$.
Следовательно, $x$ принадлежит по крайней мере одному из множеств
$A\cap B$ или $A\cap C$. Тогда $x\in(A\cap B)\cup(A\cap C)$ и, таким
образом, $A\cap(B\cup C)\subset(A\cap B)\cup(A\cap C)$. Наоборот, если
$x\in(A\cap B)\cup(A\cap C)$, то $x\in A\cap B$ или $x\in A\cap C$.
Следо\-%
\bookpage{16}%
вательно, $x\in A$ и по крайней мере одному из множеств $B$ или $C$. Но
тогда $x\in B\cup C$, а значит, и $x\in A\cap(B\cup C)$. Таким образом,
$(A\cap B)\cup(A\cap C)\subset A\cap(B\cup C)$. Сравнивая эту формулу
с полученной ранее формулой $A\cap(B\cup C)\subset(A\cap B)\cup(A\cap C)$,
приходим к формуле~$5^\circ$.

Благодаря тому что операции $A\cap B$ и $A\cup B$ ассоциативны, при
многократном их применении можно опускать скобки, например
$((A\cap B)\cap C)\cap D\equiv A\cap B\cap C\cap D$. Применяется также
сокращенная запись:
\begin{equation}\tag{1.3}
A_1\cap\dots\cap A_n\equiv\bigcap_{i=1}^{i=n}A_i;\qquad
A_1\cup\dots\cup A_n\equiv\bigcup_{i=1}^{i=n}A_i.
\end{equation}

\subsection*{Упражнения и задачи}
\label{ex:s1}
1. Доказать тождества
\begin{equation}\tag{1.4}
\begin{gathered}
A\sd B\equiv(A\cup B)\setminus(A\cap B);\\
A\cup B\equiv(A\sd B)\sd(A\cap B);\qquad A\setminus B\equiv A\sd(A\cap B);\\
(A\cap B)\cup B\equiv B;\qquad (A\cup B)\cap B\equiv B;\\
(A\cap B)\cup(A\cap\ov B)\equiv A;\qquad (A\cup B)\cap(A\cup\ov B)\equiv A;\\
(\ov A\cap B)\cup(A\cap\ov B)\equiv(A\cup B)\cap(\ov A\cup\ov B).
\end{gathered}
\end{equation}
(Для доказательства использовать формулы (1.1), (1.2) и свойства
$1^\circ$--$19^\circ$.)

2. Показать, что симметрическая разность "--- это коммутативная
и ассоциативная операции.

\textsc{Решение.} В соответствии с формулой (1.2) и свойствами
$1^\circ$--$19^\circ$ получаем
\begin{align*}
A\sd(B\sd C)&\equiv[A\cap(\ov{B\sd C})]\cup[\ov A\cap(B\sd C)]\equiv\\
&\equiv\{A\cap\ov{[(B\cap\ov C)\cup(\ov B\cap C)]}\}\cup\{\ov A\cap[(B\cap\ov C)\cup(\ov B\cap C)]\}\equiv\\
&\equiv\{A\cap[\ov{(B\cap\ov C)}\cap\ov{(\ov B\cap C)}]\}\cup\{\ov A\cap[(B\cap\ov C)\cup(\ov B\cap C)]\}\equiv\\
&\equiv\{A\cap[(\ov B\cup C)\cap(B\cup\ov C)]\}\cup\{(\ov A\cap B\cap\ov C)\cup(\ov A\cap\ov B\cap C)\}\equiv\\
&\equiv\{[A\cap(\ov B\cup C)]\cap[A\cap(B\cup\ov C)]\}\cup(\ov A\cap B\cap\ov C)\cup{}\\
&\qquad\cup(\ov A\cap\ov B\cap C)\equiv\{[(A\cap\ov B)\cup(A\cap C)]\cap[(A\cap B)\cup(A\cap\ov C)]\}\cup{}\\
&\qquad\cup(\ov A\cap B\cap\ov C)\cup(\ov A\cap\ov B\cap C)\equiv\\
&\equiv\{[(A\cap\ov B)\cap(A\cap B)]\cup[(A\cap\ov B)\cap(A\cap\ov C)]\cup[(A\cap C)\cap(A\cap B)]\cup{}\\
&\qquad\cup[(A\cap C)\cap(A\cap\ov C)]\}\cup(\ov A\cap B\cap\ov C)\cup(\ov A\cap\ov B\cap C)\equiv\\
&\equiv\{[(A\cap A)\cap(B\cap\ov B)]\cup[(A\cap A)\cap(\ov B\cap\ov C)]\cup[(A\cap A)\cap(C\cap B)]\cup{}\\
&\qquad\cup[(A\cap A)\cap(C\cap\ov C)]\}\cup(\ov A\cap B\cap\ov C)\cup(\ov A\cap\ov B\cap C)\equiv\\
&\equiv(A\cap\emptyset)\cup(A\cap\ov B\cap\ov C)\cup(A\cap B\cap C)\cup(A\cap\emptyset)\cup{}\\
&\qquad\cup(\ov A\cap B\cap\ov C)\cup(\ov A\cap\ov B\cap C)\equiv\emptyset\cup(A\cap\ov B\cap\ov C)\cup{}\\
&\qquad\cup(A\cap B\cap C)\cup\emptyset\cup(\ov A\cap B\cap\ov C)\cup(\ov A\cap\ov B\cap C)\equiv\\
&\equiv(A\cap\ov B\cap\ov C)\cup(A\cap B\cap C)\cup(\ov A\cap B\cap\ov C)\cup(\ov A\cap\ov B\cap C).
\end{align*}
\bookpage{17}Поскольку последнее выражение инвариантно относительно циклической
перестановки букв (т.\,е. при замене $A$, $B$, $C$ на $C$, $A$, $B$
соответственно), то оно равно также $C\sd(A\sd B)$. Таким образом,
ассоциативность операции $A\sd B$ доказана. Ее коммутативность
очевидна.\fixed{в оригинале в восьмой строке выкладки стоит скобка
<<$([A\cap A)$>> вместо <<$[(A\cap A)$>>.}

3. Пусть $I=\R$. Показать, что равенство
$(A\cap B)\cup C=C\cup(B\cap C)$ справедливо, если $A=[a,b]$, $B=[c,d]$,
$C=[e,f]$ и $e\le\max(a,c)$; $f\ge\min(b,d)$, но оно не является
тождеством.

4. Показать, что правила де Моргана $10^\circ$ и $11^\circ$ верны для
$n$ ($n\ge2$) компонент, т.\,е.
\begin{equation}\tag{1.5}
\ov{\bigcap_{i=1}^{i=n}A_i}\equiv\bigcup_{i=1}^{i=n}\ov{A_i};\qquad
\ov{\bigcup_{i=1}^{i=n}A_i}\equiv\bigcap_{i=1}^{i=n}\ov{A_i}.
\end{equation}

%======================================================================
\begin{solution}[1 (\S\,1)]
Будем пользоваться только (1.1), (1.2) и свойствами $1^\circ$--$19^\circ$.
Полезно заранее заметить два следствия.

(а) \emph{Поглощение:} $(A\cap B)\cup B\equiv B$ и $(A\cup B)\cap B\equiv B$.
Действительно,
$(A\cap B)\cup B\overset{14^\circ}{\equiv}(A\cap B)\cup(I\cap B)
\overset{5^\circ}{\equiv}(A\cup I)\cap B\overset{15^\circ}{\equiv}I\cap B\equiv B$;
второе тождество двойственно первому (те же шаги с $6^\circ$, $17^\circ$,
$16^\circ$):
$(A\cup B)\cap B\equiv(A\cup B)\cap(\emptyset\cup B)\equiv(A\cap\emptyset)\cup B\equiv\emptyset\cup B\equiv B$.
Это четвертое и пятое тождества (1.4).

(б) \emph{Склеивание:}
$(A\cap B)\cup(A\cap\ov B)\overset{5^\circ}{\equiv}A\cap(B\cup\ov B)
\overset{13^\circ}{\equiv}A\cap I\equiv A$;
$(A\cup B)\cap(A\cup\ov B)\overset{6^\circ}{\equiv}A\cup(B\cap\ov B)
\overset{12^\circ}{\equiv}A\cup\emptyset\equiv A$. Это шестое и седьмое тождества.

\smallskip
\emph{Последнее тождество.} Раскроем скобки в правой части по $5^\circ$:
\[
(A\cup B)\cap(\ov A\cup\ov B)\equiv(A\cap\ov A)\cup(A\cap\ov B)\cup(B\cap\ov A)\cup(B\cap\ov B)
\equiv(A\cap\ov B)\cup(\ov A\cap B),
\]
так как $A\cap\ov A\equiv B\cap\ov B\equiv\emptyset$. Попутно получена
полезная формула $A\sd B\equiv(A\cup B)\cap(\ov A\cup\ov B)$.

\emph{Первое тождество.} По (1.1) и $10^\circ$:
$(A\cup B)\setminus(A\cap B)\equiv(A\cup B)\cap\ov{A\cap B}\equiv(A\cup B)\cap(\ov A\cup\ov B)$,
а это, как только что показано, $A\sd B$.

\emph{Третье тождество.} По (1.2), $10^\circ$ и $9^\circ$:
\begin{multline*}
A\sd(A\cap B)\equiv\bigl(A\cap\ov{A\cap B}\bigr)\cup\bigl(\ov A\cap A\cap B\bigr)
\equiv\bigl(A\cap(\ov A\cup\ov B)\bigr)\cup\emptyset\equiv{}\\
\equiv(A\cap\ov A)\cup(A\cap\ov B)\equiv A\cap\ov B\equiv A\setminus B.
\end{multline*}

\emph{Второе тождество.} Обозначим $D=A\sd B=(A\cap\ov B)\cup(\ov A\cap B)$,
$P=A\cap B$. Заметим, что $D\cap P\equiv\emptyset$ (каждый член $D$
содержит $\ov A$ или $\ov B$). Тогда $D\cap\ov P\equiv D$ (так как
$D\equiv(D\cap P)\cup(D\cap\ov P)$ по склеиванию) и аналогично
$P\cap\ov D\equiv P$. Поэтому
\[
(A\sd B)\sd(A\cap B)\equiv(D\cap\ov P)\cup(\ov D\cap P)\equiv D\cup P
\equiv(A\cap\ov B)\cup(\ov A\cap B)\cup(A\cap B).
\]
Склеивая $(A\cap\ov B)\cup(A\cap B)\equiv A$ и $(\ov A\cap B)\cup(A\cap B)\equiv B$
(член $A\cap B$ можно повторить по $8^\circ$), получаем $A\cup B$.
\end{solution}

\begin{solution}[2 (\S\,1): дополнение]
Решение в книге приводит $A\sd(B\sd C)$ к виду
$K_{100}\cup K_{111}\cup K_{010}\cup K_{001}$, где
$K_{\alpha\beta\gamma}=A^\alpha\cap B^\beta\cap C^\gamma$ ($X^1=X$,
$X^0=\ov X$). Это ровно те <<кусочки>> диаграммы Эйлера, в которых точка
принадлежит \emph{нечетному} числу множеств. Отсюда сразу видна и
коммутативность, и ассоциативность: для любых $A_1,\dots,A_n$ множество
$A_1\sd\dots\sd A_n$ (при любой расстановке скобок) "--- это множество
точек, принадлежащих нечетному числу из $A_i$. (Индукция: $x\in X\sd Y$
$\iff$ $x$ лежит ровно в одном из $X$, $Y$, т.\,е. число вхождений в $X\sd Y$
имеет ту же четность, что сумма чисел вхождений.) Коммутативность $A\sd B\equiv B\sd A$
следует из (1.2) и $1^\circ$, $2^\circ$.
\end{solution}

\begin{solution}[3 (\S\,1)]
По поглощению $C\cup(B\cap C)\equiv C$, поэтому равенство
$(A\cap B)\cup C=C\cup(B\cap C)$ равносильно $(A\cap B)\cup C=C$, т.\,е.
\[
A\cap B\subset C .
\]
Для отрезков $A\cap B=[\max(a,c),\ \min(b,d)]$ (пусто, если
$\max(a,c)>\min(b,d)$). Если $e\le\max(a,c)$ и $f\ge\min(b,d)$, то
$[\max(a,c),\min(b,d)]\subset[e,f]$, и равенство выполнено (если же
$A\cap B=\emptyset$, оно тем более верно).

Тождеством равенство не является: при $A=B=[0,1]$, $C=[2,3]$ слева
$[0,1]\cup[2,3]$, справа $[2,3]$. Вообще, равенство верно в точности при
$A\cap B\subset C$; условия задачи "--- достаточные (а при
$A\cap B\neq\emptyset$ и необходимые).
\end{solution}

\begin{solution}[4 (\S\,1)]
Индукция по $n$. При $n=2$ это $10^\circ$, $11^\circ$. Пусть (1.5) верно
для $n$. Тогда по ассоциативности, $10^\circ$ и предположению индукции
\[
\ov{\bigcap_{i=1}^{n+1}A_i}\equiv\ov{\Bigl(\bigcap_{i=1}^{n}A_i\Bigr)\cap A_{n+1}}
\equiv\ov{\bigcap_{i=1}^{n}A_i}\cup\ov{A_{n+1}}
\equiv\Bigl(\bigcup_{i=1}^{n}\ov{A_i}\Bigr)\cup\ov{A_{n+1}}\equiv\bigcup_{i=1}^{n+1}\ov{A_i}.
\]
Второе тождество получается так же из $11^\circ$ (или из первого,
примененного к $\ov{A_i}$, с последующим взятием дополнения и законом
$9^\circ$). Можно рассуждать и непосредственно: $x\notin\bigcap A_i$
$\iff$ $x\notin A_i$ хотя бы для одного $i$ $\iff$ $x\in\bigcup\ov{A_i}$.
\end{solution}

\section{Диаграммы Эйлера}

Логические операции над множествами допускают простую геометрическую
интерпретацию. В качестве унитарного множества $I$ возьмем плоскость, на
которой будем изображать множества в виде некоторых областей. Тогда
множество $A\cap B$ будет представлять собой область (заштрихованную),
показанную на рис.~1,~\emph{а}. Области (заштрихованные), соответствующие
операциям $A\cup B$, $A\setminus B$, $A\sd B$ и $\ov A$, показаны на рис.~1,
\emph{б}, \emph{в}, \emph{г} и \emph{д} соответственно.

\begin{figure}[h]
\centering
\tikzset{hat/.style={pattern=north east lines,pattern color=black!70}}
\def\ca{(0,0) circle (0.8)}
\def\cb{(0.7,-0.5) circle (0.8)}
\begin{tikzpicture}[scale=0.9]
 \begin{scope}\clip\ca; \fill[hat]\cb; \end{scope}
 \draw\ca node[left=6pt]{$A$}; \draw\cb node[right=4pt]{$B$};
 \node at (0.35,-1.6){\emph{а}}; \node at (0.35,-1.95){\small $A\cap B$};
\end{tikzpicture}\quad
\begin{tikzpicture}[scale=0.9]
 \fill[hat]\ca; \fill[hat]\cb;
 \draw\ca node[left=6pt]{$A$}; \draw\cb node[right=4pt]{$B$};
 \node at (0.35,-1.6){\emph{б}}; \node at (0.35,-1.95){\small $A\cup B$};
\end{tikzpicture}\quad
\begin{tikzpicture}[scale=0.9]
 \begin{scope}[even odd rule]\clip\cb (-2,-2) rectangle (2,2); \fill[hat]\ca; \end{scope}
 \draw\ca node[left=6pt]{$A$}; \draw\cb node[right=4pt]{$B$};
 \node at (0.35,-1.6){\emph{в}}; \node at (0.35,-1.95){\small $A\setminus B$};
\end{tikzpicture}\quad
\begin{tikzpicture}[scale=0.9]
 \fill[hat,even odd rule]\ca \cb;
 \draw\ca node[left=6pt]{$A$}; \draw\cb node[right=4pt]{$B$};
 \node at (0.35,-1.6){\emph{г}}; \node at (0.35,-1.95){\small $A\sd B$};
\end{tikzpicture}\quad
\begin{tikzpicture}[scale=0.9]
 \fill[hat,even odd rule](-1.2,-1.2) rectangle (1.2,1.2) (0,0) circle (0.7);
 \draw (0,0) circle (0.7) node{$A$};
 \node at (0,-1.6){\emph{д}}; \node at (0,-1.95){\small $\ov A$};
\end{tikzpicture}
\figcap{Рис. 1.}
\end{figure}

\begin{figure}[h]
\centering
\begin{tikzpicture}[scale=0.9]
 \def\a{(0,0) circle (0.9)}\def\b{(1.2,0) circle (0.9)}\def\c{(0.6,-1.0) circle (0.9)}
 \fill[pattern=north east lines,pattern color=black!70,even odd rule]\a\b\c;
 \begin{scope}\clip\a;\clip\b;\fill[pattern=north east lines,pattern color=black!70]\c;\end{scope}
 \draw\a node[left=6pt]{$A$};\draw\b node[right=6pt]{$B$};\draw\c node[below=6pt]{$C$};
\end{tikzpicture}
\figcap{Рис. 2.}
\end{figure}

\bookpage{18}%
Применение этих диаграмм, названных диаграммами Эйлера, во многих случаях
позволяет существенно упростить решение задач.

Рассмотрим, например, упражнение~2 \S\,1. Формальное доказательство
ассоциативности симметрической разности $A\sd B$ потребовало выполнения ряда
тождественных преобразований. Построим диаграмму Эйлера для выражения
$A\sd(B\sd C)$ (рис.~2). Из диаграммы видно, что область $A\sd(B\sd C)$
симметрична относительно множеств $A$, $B$, $C$, и такую же диаграмму имеет
множество $(A\sd B)\sd C$.

\begin{figure}[h]
\centering
\begin{tikzpicture}[scale=0.85]
 \draw (0,0) circle (0.8); \draw (0.8,-0.6) circle (0.8);
 \node at (-0.45,0.35){$A$}; \node at (-0.25,-0.15){\small 2};
 \node at (0.45,-0.3){\small 1}; \node at (1.05,-0.95){\small 3};
 \node at (1.25,-0.4){$B$}; \node at (1.7,0.7){\small 4};
 \node at (0.4,-1.8){\emph{а}};
\end{tikzpicture}\qquad\qquad
\begin{tikzpicture}[scale=0.85]
 \draw (0,0) circle (0.9) node[left=4pt]{$A$};
 \draw (1.2,0) circle (0.9) node[right=4pt]{$B$};
 \draw (0.6,-1.0) circle (0.9); \node at (0.95,-1.4){$C$};
 \node at (0.6,0.45){\small 2}; \node at (0.6,-0.35){\small 1};
 \node at (0.05,-0.6){\small 4}; \node at (1.15,-0.6){\small 3};
 \node at (0.35,-1.45){\small 6}; \node at (-0.9,-0.9){\small 7};
 \node at (2.1,-0.6){\small 5}; \node at (2.3,-1.4){\small 8};
 \node at (0.6,-2.3){\emph{б}};
\end{tikzpicture}
\figcap{Рис. 3.}
\end{figure}

Заметим, что диаграммами Эйлера удобно пользоваться лишь при небольшом
числе рассматриваемых областей. С увеличением их числа трудность построения
диаграмм быстро возрастает. Дело в том, что граница каждой следующей
области, изображаемой на диаграмме, должна обязательно пересекаться
с каждой из логически возможных подобластей, на которые предыдущие заданные
области разбивают плоскость. (В противном случае некоторые логически
возможные ситуации могут быть пропущены.) Так, одной области $A$
соответствуют две подобласти плоскости $\R^2$: сама область $A$ и ее
внешность $\ov A$. Двум областям $A$ и $B$ соответствуют уже четыре
подобласти (рис.~3,~\emph{а}), а трем подобластям "--- восемь
(рис.~3,~\emph{б}). Такие области будем называть \emph{конституентными
областями}. Очевидно, подобрать четвертую область $D$, граница которой
проходила бы через внутренние точки каждой из подобластей, уже труднее
(рис.~4). Заметим, что если, например, вместо рис.~3,~\emph{б} принять
расположение областей такое, как показано на рис.~5, то логической формуле
$(A\cap B)\cap C$ соответствовало бы пустое множество $\emptyset$, что для
других множеств $A$, $B$ и $C$, вообще, не верно.

\begin{figure}[h]
\centering
\begin{tikzpicture}[scale=0.8]
 \draw (0,0) circle (0.9) node[left=4pt]{$A$};
 \draw (1.2,0) circle (0.9) node[right=4pt]{$B$};
 \draw (0.6,-1.0) circle (0.9) node[right=2pt]{$C$};
 \draw[dashed,rotate around={-60:(0.4,-0.6)}] (0.4,-0.6) ellipse (1.9 and 0.5);
 \node at (-0.75,-1.9){$D$};
\end{tikzpicture}
\figcap{Рис. 4.}
\begin{tikzpicture}[scale=0.8]
 \draw (0,0) circle (0.8) node{$A$};
 \draw (1.3,0) circle (0.8) node{$B$};
 \draw (2.6,0) circle (0.8) node{$C$};
\end{tikzpicture}
\figcap{Рис. 5.}
\end{figure}

\subsection*{Упражнения и задачи}
1. Используя диаграммы Эйлера, доказать тождества\fixed{в оригинале последнее тождество напечатано как
$A\cup[(B\setminus\ov A)\sd C]\equiv(A\cap B)\sd C$, которое тождеством
не является (при $A=I$, $B=C=\emptyset$ левая часть равна $I$, правая "---
$\emptyset$). Восстановленный вариант и разбор "--- в решении ниже.}
\[
\begin{gathered}
(A\cup C)\cap B\equiv(A\cap B)\cup(\ov A\cap B\cap C);\\
[A\cap(B\cup C)]\cup(\ov A\cap C)\equiv C\cup(A\cap B);\\
(A\cap C)\cup(B\cap C)\cup(C\setminus A)\equiv C;\qquad A\sd B\equiv(A\cup B)\setminus(A\cap B);\\
[A\cap(B\sd C)]\cup(B\sd C)\equiv B\sd C;\\
A\cap[(B\setminus\ov A)\sd C]\equiv(A\cap B)\sd(A\cap C).
\end{gathered}
\]


\bookpage{19}%
2. Выпуклые области $A,B,C\subset\R^2$ разбивают плоскость на восемь частей,
различающихся состояниями принадлежности их точек множествам $A$, $B$ и $C$.
Может ли добавление выпуклой области $D$ привести к разбиению $\R^2$ на 16
подобластей указанного типа? Какое максимальное число таких подобластей
могут образовать $n$ выпуклых областей в $\R^2$?

3. Существуют ли (для всякого $n$) конгруэнтные (т.\,е. получающиеся одна из
другой переносом и поворотом) области $A_1,\dots,A_n\subset\R^2$,
разбивающие плоскость на максимально возможное ($2^n$) число частей, которые
различаются между собой состояниями принадлежности их точек областям
$A_1,\dots,A_n$. (Заметим, что для $n\le3$ такими областями, например,
являются круги одинакового радиуса.)

%======================================================================
\begin{solution}[1 (\S\,2)]
Нумеруем конституентные области трех кругов, как на рис.~3,~\emph{б},
и записываем, какие из них образуют левую и правую части. Удобнее
обозначать область тройкой $\alpha\beta\gamma$ состояний принадлежности
к $A$, $B$, $C$ ($1$ "--- внутри, $0$ "--- вне); тогда
\[
1=111,\ 2=110,\ 3=011,\ 4=101,\ 5=010,\ 6=001,\ 7=100,\ 8=000 .
\]
\begin{center}\small
\begin{tabular}{lll}
\toprule
тождество & левая часть & правая часть\\
\midrule
$(A\cup C)\cap B\equiv(A\cap B)\cup(\ov A\cap B\cap C)$ & $\{1,2,3\}$ & $\{1,2\}\cup\{3\}$\\
$[A\cap(B\cup C)]\cup(\ov A\cap C)\equiv C\cup(A\cap B)$ & $\{1,2,4\}\cup\{3,6\}$ & $\{1,3,4,6\}\cup\{1,2\}$\\
$(A\cap C)\cup(B\cap C)\cup(C\setminus A)\equiv C$ & $\{1,4\}\cup\{1,3\}\cup\{3,6\}$ & $\{1,3,4,6\}$\\
$A\sd B\equiv(A\cup B)\setminus(A\cap B)$ & $\{4,7,3,5\}$ & $\{1,\dots,5,7\}\setminus\{1,2\}$\\
$[A\cap(B\sd C)]\cup(B\sd C)\equiv B\sd C$ & $\{2,4\}\cup\{2,4,5,6\}$ & $\{2,4,5,6\}$\\
\bottomrule
\end{tabular}
\end{center}
Во всех строках множества номеров совпадают, а так как три круга
в положении рис.~3,~\emph{б} реализуют все восемь логически возможных
состояний, совпадение областей на диаграмме и означает тождество
(см. также упр.~2 \S\,3 и решение упр.~7 \S\,3, где это рассуждение
обосновано строго). Четвертое и пятое тождества видны уже на диаграмме
двух кругов (пятое "--- это поглощение $(X\cap Y)\cup Y\equiv Y$).

\emph{Последнее тождество.} Поскольку $B\setminus\ov A\equiv B\cap A$,
напечатанное в книге равенство
$A\cup[(B\setminus\ov A)\sd C]\equiv(A\cap B)\sd C$ имеет левую часть
$A\cup\{1,2\}\sd\{1,3,4,6\}=\{1,2,4,7\}\cup\{2,3,4,6\}=\{1,2,3,4,6,7\}$,
а правую $\{2,3,4,6\}$: они различаются областями $1$ и $7$ (например, при
$A=I$, $B=C=\emptyset$). Значит, в книге опечатка. Естественный
верный вариант "--- с пересечением вместо объединения:
\[
A\cap[(B\setminus\ov A)\sd C]\equiv(A\cap B)\sd(A\cap C):
\]
слева $\{1,2,4,7\}\cap\{2,3,4,6\}=\{2,4\}$, справа
$\{1,2\}\sd\{1,4\}=\{2,4\}$. Это тождество выражает
дистрибутивность $\cap$ относительно $\sd$:
$A\cap(X\sd Y)\equiv(A\cap X)\sd(A\cap Y)$. (Верно также
$A\cap[(B\setminus\ov A)\sd C]\equiv(A\cap B)\sd(A\cap C)\equiv A\cap(B\sd C)$.)
\end{solution}

\begin{solution}[2 (\S\,2)]
\emph{Ответ: да; $n$ выпуклых областей могут дать все $2^n$
логически возможных частей.}

Для $n=4$ это показывает, например, известная диаграмма Венна из четырех
равных эллипсов (эллипс выпуклый). Общая конструкция для любого $n$:
возьмем на окружности $2^n$ различных точек $q_\alpha$, занумерованных
наборами $\alpha=(\alpha_1,\dots,\alpha_n)\in\{0,1\}^n$, и положим
\[
A_i=\operatorname{conv}\{q_\alpha:\ \alpha_i=1\}
\]
(выпуклая оболочка; чтобы получить область с внутренними точками, ее можно
слегка <<раздуть>>, взяв $\varepsilon$-окрестность при малом
$\varepsilon$ "--- выпуклость сохраняется). Точки на окружности находятся
в выпуклом положении: ни одна из них не лежит в выпуклой оболочке
остальных и даже отстоит от нее на положительное расстояние. Поэтому
$q_\alpha\in A_i\iff\alpha_i=1$, т.\,е. точка $q_\alpha$ принадлежит
конституенте $A_1^{\alpha_1}\cap\dots\cap A_n^{\alpha_n}$, и все $2^n$
конституент непусты (а вместе с малыми окрестностями точек $q_\alpha$,
лежащими вне всех $A_i$ с $\alpha_i=0$, "--- имеют внутренние точки).

Итак, максимальное число частей, различающихся состояниями
принадлежности, равно $2^n$ "--- больше быть не может, так как наборов
состояний всего $2^n$. (Связными эти части, вообще говоря, не будут; если
же считать связные компоненты, то их число для выпуклых областей не
ограничено никакой функцией от $n$: границы двух выпуклых многоугольников
могут пересекаться в сколь угодно большом числе точек.)
\end{solution}

\begin{solution}[3 (\S\,2)]
\emph{Ответ: да, для всякого $n$; можно даже обойтись одними
параллельными переносами.}

Выберем $n$ векторов $v_1,\dots,v_n$ и $2^n$ точек $x_\alpha$,
$\alpha\in\{0,1\}^n$, <<в общем положении>>, так чтобы все $n2^n$ точек
$x_\alpha-v_i$ были попарно различны (этого можно добиться сколь угодно
малым шевелением). Разобьем их на два конечных множества
\[
P_+=\{x_\alpha-v_i:\ \alpha_i=1\},\qquad P_-=\{x_\alpha-v_i:\ \alpha_i=0\}.
\]
Пусть $S$ "--- связная открытая область, содержащая все точки $P_+$ и не
содержащая ни одной точки $P_-$ (например, объединение малых кругов вокруг
точек $P_+$, соединенных тонкими <<коридорами>>, обходящими конечное
множество $P_-$, "--- плоскость без конечного множества точек связна).
Положим $A_i=S+v_i$. Тогда
\[
x_\alpha\in A_i\iff x_\alpha-v_i\in S\iff\alpha_i=1,
\]
т.\,е. $x_\alpha$ лежит в конституенте с номером $\alpha$, и все $2^n$
конституент непусты (и открыты в окрестности $x_\alpha$, если круги
выбрать малыми). Области $A_i$ конгруэнтны "--- они получаются друг из
друга переносом.

Для $n\le3$ подходят равные круги (рис.~3,~\emph{б}). Для $n\ge4$
равные (и вообще любые) \emph{круги} уже не годятся: $n$ окружностей
разбивают плоскость не более чем на $n^2-n+2$ связных частей (каждая новая,
$k$-я, окружность пересекает прежние не более чем в $2(k-1)$ точках
и добавляет не более $2(k-1)$ частей). Каждая непустая конституента
содержит хотя бы одну такую часть, поэтому непустых конституент не больше
$n^2-n+2$, а $n^2-n+2<2^n$ при $n\ge4$. Поэтому при $n\ge4$ приходится
брать области более сложной формы, как в построении выше. Для простых $n$
существуют даже симметричные диаграммы Венна из $n$ конгруэнтных кривых,
переходящих друг в друга поворотом (Григгс, Киллиан, Сэвидж, 2004).
\end{solution}

\section{Принцип двойственности}

Пусть $F(A_1,\dots,A_n)$ "--- некоторая логическая формула, написанная
с помощью символов (не обязательно всех) операций пересечения, соединения
и дополнения, пустого $\emptyset$ и унитарного $I$ множеств. (Если
в $F(A_1,\dots,A_n)$ присутствуют символы разности $A\setminus B$
и симметрической разности $A\sd B$, то их можно с помощью формул (1.1)
и (1.2) исключить.)

Формулу $F^*(A_1,\dots,A_n)$ будем называть \emph{двойственной} формуле
$F(A_1,\dots,A_n)$, если она может быть получена из последней путем
формальной замены (<<преобразования двойственности>>) символов ($\cap$ на
$\cup$, $\cup$ на $\cap$, $\emptyset$ на $I$, $I$ на $\emptyset$) и,
возможно, последующих тождественных преобразований. Примерами двойственных
формул являются
\[
\begin{gathered}
(A\cap B)\cup(\ov A\cap I)\quad\text{и}\quad(A\cup B)\cap(\ov A\cup\emptyset);\\
\ov{(\ov A\cup B)\cap C}\quad\text{и}\quad\ov{(\ov A\cap B)\cup C};\\
(A\cap B)\cup A\quad\text{и}\quad A.
\end{gathered}
\]
Если $F(A_1,\dots,A_n)\equiv F^*(A_1,\dots,A_n)$ (т.\,е. двойственные
формулы тождественны), формулу $F(A_1,\dots,A_n)$ будем называть
\emph{самодвойственной}. Примером самодвойственной формулы является
\begin{equation}\tag{1.6}
(A\cap B)\cup(A\cap C)\cup(B\cap C).
\end{equation}
Действительно,\fixed{в оригинале четвертый сомножитель после раскрытия скобок напечатан как
$(B\cup C\cup C)$; это сомножитель $(B\cup C)\cup B$, т.\,е.
$(B\cup C\cup B)$ (на результат не влияет).}
\begin{align*}
&(A\cap B)\cup(A\cap C)\cup(B\cap C)\equiv[(A\cup A)\cap(A\cup C)\cap(B\cup A)\cap(B\cup C)]\cup(B\cap C)\equiv\\
&\equiv[A\cap(A\cup C)\cap(A\cup B)\cap(B\cup C)]\cup(B\cap C)\equiv(A\cup B)\cap(A\cup B\cup C)\cap{}\\
&\cap(A\cup B\cup B)\cap(B\cup C\cup B)\cap(A\cup C)\cap(A\cup C\cup C)\cap{}\\
&\cap(A\cup B\cup C)\cap(B\cup C\cup C)\equiv(A\cup B)\cap(A\cup B\cup C)\cap{}\\
&\cap(A\cup B)\cap(B\cup C)\cap(A\cup C)\cap(A\cup C)\cap(A\cup B\cup C)\cap{}\\
&\cap(B\cup C)\equiv(A\cup B)\cap(A\cup B\cup C)\cap(B\cup C)\cap(A\cup C)\equiv{}\\
&\equiv(A\cup B)\cap(B\cup C)\cap(A\cup C).
\end{align*}

\bookpage{20}%
Очевидно, также, что если $F(A_1,\dots,A_n)=A_k$ ($1\le k\le n$), то
$F(A_1,\dots,A_n)$ "--- самодвойственная формула.

Тождества в алгебре множеств можно рассматривать как теоремы, путь
к доказательствам которых лежит через тождественные преобразования.
(В исчислении высказываний этим преобразованиям соответствуют элементарные
рассуждения формальной логики [24, 84].) Пусть
\begin{equation}\tag{1.7}
F(A_1,\dots,A_n)\equiv\Phi(A_1,\dots,A_n)
\end{equation}
"--- тождество, записанное с помощью символов $\cap$, $\cup$,
${}^{\overline{\phantom{a}}}$, $\emptyset$ и $I$. Нетрудно заметить, что
если в тождествах $1^\circ$--$19^\circ$ \S\,1 произведем упомянутое выше
преобразование двойственности (т.\,е. заменим $\cap$ на $\cup$, $\cup$ на
$\cap$, $\emptyset$ на $I$ и $I$ на $\emptyset$), то получим ту же систему
тождеств $1^\circ$--$19^\circ$. Отсюда следует, что при этом
преобразовании тождество (1.7) также сохраняется:
\begin{equation}\tag{1.8}
F^*(A_1,\dots,A_n)\equiv\Phi^*(A_1,\dots,A_n).
\end{equation}
(Таким образом, каждой теореме в алгебре множеств соответствует
двойственная теорема.)

Заметим, что если множества, определяемые формулами $F(A_1,\dots,A_n)$
и $\Phi(A_1,\dots,A_n)$, совпадают лишь при некотором фиксированном выборе
множеств $A_1,\dots,A_n$ (т.\,е. (1.7) является равенством, а не
тождеством), то двойственные им формулы будут, вообще, определять
различные множества.

\tikzset{hat/.style={pattern=north east lines,pattern color=black!70}}
\begin{figure}[h]
\centering
\begin{tikzpicture}[scale=0.75]
 \fill[hat] (-1,0) arc (180:0:1) -- cycle;
 \fill[hat] (2,0) arc (0:-180:2) -- cycle;
 \draw (-1,0) arc (180:0:1);
 \draw (2,0) arc (0:-180:2);
 \draw[dashed] (-2,0) -- (-1,0) (1,0) -- (2,0);
 \draw (-1,0) -- (1,0);
 \draw[->] (-2.6,0) -- (2.8,0) node[below]{$x_1$};
 \draw[->] (0,-2.4) -- (0,1.6) node[left]{$x_2$};
 \node[below right] at (0,0){\small 0}; \node[above] at (1.15,0){\small 1}; \node[above] at (2.15,0){\small 2};
 \node at (0.9,-1.0){$\Omega$};
\end{tikzpicture}
\figcap{Рис. 6.}
\end{figure}

{\small
\noindent\textbf{Пример 1.} Пусть $A$, $B$ и $C$ "--- области на
плоскости $x_1Ox_2$, определяемые неравенствами
$A\equiv(4-x_1^2-x_2^2\ge0)$ (круг радиуса два с центром в начале
координат); $B\equiv(x_2\ge0)$ (верхняя полуплоскость);
$C\equiv(1-x_1^2-x_2^2\ge0)$ (круг радиуса один с центром в начале
координат). Легко проверить, что множества
$M\equiv(C\cap B)\cup(A\cap\ov B)$ и $N\equiv(C\cup\ov B)\cap A$ совпадают
и определяют область $\Omega$ (заштрихованная часть на рис.~6;
штриховыми линиями обозначены участки границы области $\Omega$, не
принадлежащие $\Omega$). Поэтому $(C\cap B)\cup(A\cap\ov B)=(C\cup\ov B)\cap A$.
Рассмотрим двойственную формулу $(C\cup B)\cap(A\cup\ov B)=(C\cap\ov B)\cup A$.
Нетрудно убедиться в том, что левая часть этой формулы соответствует
область, приведенная на рис.~7,~\emph{а}, а правой "--- на рис.~7,~\emph{б}.
Таким образом, в этом случае двойственное равенство не выполняется.

\smallskip
\noindent\textbf{Пример 2.} Рассмотрим множество
\[
M=(A_1\cap\ov A_2\cap\ov A_3)\cup(A_1\cap A_2)\cup
\bigl[\,\ov{\ov{A_1\cap A_3\cap I}\cap\ov{A_2\cap A_3}}\,\bigr].
\]
Применяя формулы $1^\circ$--$19^\circ$ \S\,1, последовательно получаем
\begin{align*}
M&\equiv(A_1\cap\ov A_2\cap\ov A_3)\cup(A_1\cap A_2)\cup(A_1\cap A_3\cap I)\cup(A_2\cap A_3)\equiv\\
&\equiv(A_1\cap\ov A_2\cap\ov A_3)\cup(A_1\cap A_2)\cup(A_1\cap A_3)\cup(A_2\cap A_3)\equiv\\
&\equiv(A_1\cap\ov A_2\cap\ov A_3)\cup[A_1\cap(A_2\cup A_3)]\cup(A_2\cap A_3)\equiv\\
&\equiv[A_1\cap(\ov A_2\cap\ov A_3)]\cup[A_1\cap(A_2\cup A_3)]\cup(A_2\cap A_3)\equiv\\
&\equiv\{A_1\cap[(\ov A_2\cap\ov A_3)\cup(A_2\cup A_3)]\}\cup(A_2\cap A_3)\equiv\\
&\equiv\{A_1\cap\{[(\ov A_2\cap\ov A_3)\cup A_2]\cup A_3\}\}\cup(A_2\cap A_3)\equiv\\
&\equiv\{A_1\cap\{[(\ov A_2\cup A_2)\cap(\ov A_3\cup A_2)]\cup A_3\}\}\cup(A_2\cap A_3)\equiv\\
&\equiv\{A_1\cap\{[I\cap(\ov A_3\cup A_2)]\cup A_3\}\}\cup(A_2\cap A_3)\equiv\\
&\equiv[A_1\cap(A_2\cup A_3\cup\ov A_3)]\cup(A_2\cap A_3)\equiv\\
&\equiv[A_1\cap(A_2\cup I)]\cup(A_2\cap A_3)\equiv(A_1\cap I)\cup(A_2\cap A_3)\equiv A_1\cup(A_2\cap A_3).
\end{align*}
}

\bookpage{21}\begin{figure}[h]
\centering
\begin{tikzpicture}[scale=0.7]
 \fill[hat] (-2,0) arc (180:0:2) -- cycle;
 \fill[hat] (1,0) arc (0:-180:1) -- cycle;
 \draw (-2,0) arc (180:0:2); \draw (1,0) arc (0:-180:1);
 \draw[dashed] (2,0) arc (0:-180:2);
 \draw (-2,0)--(2,0);
 \draw[->] (-2.6,0) -- (2.8,0) node[below]{$x_1$};
 \draw[->] (0,-2.4) -- (0,2.5) node[left]{$x_2$};
 \node[below] at (1.2,0){\small 1}; \node[below] at (2.2,0){\small 2};
 \node at (0,-2.8){\emph{а}};
\end{tikzpicture}\qquad\qquad
\begin{tikzpicture}[scale=0.7]
 \fill[hat] (0,0) circle (2); \draw (0,0) circle (2);
 \draw[->] (-2.6,0) -- (2.8,0) node[below]{$x_1$};
 \draw[->] (0,-2.4) -- (0,2.5) node[left]{$x_2$};
 \node at (0,-2.8){\emph{б}};
\end{tikzpicture}
\figcap{Рис. 7.}
\end{figure}

{\small\noindent Таким образом, справедливо тождество
\[
(A_1\cap\ov A_2\cap\ov A_3)\cup(A_1\cap A_2)\cup
\bigl[\,\ov{\ov{A_1\cap A_3\cap I}\cap\ov{A_2\cap A_3}}\,\bigr]\equiv A_1\cup(A_2\cap A_3).
\]
Следовательно, выполнимо и двойственное тождество
\[
(A_1\cup\ov A_2\cup\ov A_3)\cap(A_1\cup A_2)\cap
\bigl[\,\ov{\ov{A_1\cup A_3\cup\emptyset}\cup\ov{A_2\cup A_3}}\,\bigr]\equiv A_1\cap(A_2\cup A_3).
\]\par}

\subsection*{Упражнения и задачи}
1. Существует ли самодвойственная двухместная операция $F(A,B)$, не равная
тождественно $A$ или $B$?

2. Описать множество всех самодвойственных формул вида $F(A_1,\dots,A_n)$
при любом $n\ge2$.

3. Описать множество всех логических формул, для которых
$F\cap F^*=\emptyset$. Будет ли при этом $F\cup F^*=I$?

4.\fixed{в оригинале номер упражнения <<4.>> пропущен.} Будем
предполагать, что $F(A_1,\dots,A_n)$ "--- \emph{строго $n$-местная}
логическая операция, если не существует таких целых чисел
$k_1,\dots,k_s\in[1,n]$ ($s<n$), что
$F(A_1,\dots,A_n)=\Phi(A_{k_1},\dots,A_{k_s})$, где $\Phi$ "--- $s$-местная
логическая формула. Для каких $n$ существуют строго $n$-местные
самодвойственные логические формулы? Простой перебор показывает, что для
$n=2$ таких формул нет. Для $n=3$ примером самодвойственной строго
трехместной формулы является (1.6).

5. Пусть $\mu(n)$ "--- количество самодвойственных $n$-местных логических
формул. Показать, что $\mu(n)$ "--- неубывающая функция.
\bookpage{22}%

6. Пусть $\tau(n)$ "--- количество самодвойственных строго $n$-местных
логических формул. Показать, что $\tau(2)=0$. Для каких еще $n$
$\tau(n)=0$?

7. Пусть $\lambda\{F(A_1,\dots,A_n)\}=\lambda(F)$ "--- количество
конституентных областей в диаграмме Эйлера для логической формулы
$F(A_1,\dots,A_n)$, а $\lambda(F^*)$ "--- количество таких же областей для
двойственной формулы. Описать класс функций, для которых
$\lambda(F)=\lambda(F^*)$.

% Решения упражнений § 3 (с. 21--22).
% Основа: заметки claude/p21-*.md из проекта (см. каталог notes/).

\begin{remarkbox}{Общие инструменты для упражнений 1--7}
Во всех упражнениях формулы рассматриваются \emph{с точностью до
тождественного равенства}: две формулы, тождественно равные друг другу,
считаются одной и той же (они задают одну и ту же <<операцию над
множествами>>).

\textbf{Лемма (L) о двойственности.} Для любой формулы $F$
\begin{equation*}\tag{L}
F^*(A_1,\dots,A_n)\equiv\ov{F(\ov{A_1},\dots,\ov{A_n})}.
\end{equation*}

\emph{Доказательство} "--- полная индукция по числу $k(F)$ символов
операций $\cap$, $\cup$, ${}^{\overline{\phantom a}}$ в записи $F$.
Обозначим $\widehat F=\ov{F(\ov A)}$, где $\ov A=(\ov{A_1},\dots,\ov{A_n})$.

\emph{База} $k=0$: $F=A_i$ дает $F^*=A_i=\ov{\ov{A_i}}=\widehat F$;
$F=\emptyset$ дает $F^*=I=\ov\emptyset=\widehat F$; $F=I$ дает
$F^*=\emptyset=\ov I=\widehat F$.

\emph{Шаг.} Пусть (L) верно для всех формул с не более чем $k$ операциями
и $k(F)=k+1$. Внешняя операция $F$ "--- одна из трех.
\begin{itemize}[nosep]
\item $F=\ov G$, $k(G)=k$. Замена символов черту не трогает, поэтому
  $F^*=\ov{G^*}\equiv\ov{\ov{G(\ov A)}}=G(\ov A)$, и
  $\widehat F=\ov{\ov{G(\ov A)}}=G(\ov A)$.
\item $F=G\cap H$, $k(G)+k(H)=k$, так что $k(G),k(H)\le k$ (именно здесь
  нужна \emph{полная} индукция: подформулы могут быть намного короче).
  Тогда $F^*=G^*\cup H^*\equiv\ov{G(\ov A)}\cup\ov{H(\ov A)}
  \overset{10^\circ}{\equiv}\ov{G(\ov A)\cap H(\ov A)}=\widehat F$.
\item $F=G\cup H$ "--- так же, с правилом $11^\circ$. \qquad$\square$
\end{itemize}

\textbf{Коды точек и булева функция формулы.} Точке $x\in I$ сопоставим
\emph{код} $\alpha(x)=(\alpha_1,\dots,\alpha_n)\in\{0,1\}^n$, $\alpha_i=1$
при $x\in A_i$ и $\alpha_i=0$ иначе. Точки с одним кодом $\alpha$ образуют
\emph{конституенту}
\[
K_\alpha=A_1^{\alpha_1}\cap\dots\cap A_n^{\alpha_n},\qquad A^1=A,\ A^0=\ov A ;
\]
конституент $2^n$, они попарно не пересекаются и покрывают $I$.
Принадлежность точки множеству $F(A_1,\dots,A_n)$ зависит только от ее
кода (индукция по $k(F)$), поэтому формуле отвечает булева функция
$f\colon\{0,1\}^n\to\{0,1\}$ (таблица истинности; $\cap\to\wedge$,
$\cup\to\vee$, черта $\to$ отрицание, $\emptyset\to0$, $I\to1$) и
\begin{equation*}\tag{D}
x\in F\iff f(\alpha(x))=1,\qquad F=\bigcup_{\alpha:\,f(\alpha)=1}K_\alpha .
\end{equation*}

\emph{(i) Тождества = равенство функций.} $F\equiv G\iff f=g$. Если
$f(\alpha_0)\ne g(\alpha_0)$, возьмем $I=\{0,1\}^n$,
$A_i=\{\alpha:\alpha_i=1\}$; тогда $K_{\alpha_0}=\{\alpha_0\}$, и точка
$\alpha_0$ лежит ровно в одном из множеств $F$, $G$.

\emph{(ii) Функция двойственной формулы.} Замена $A_i\to\ov{A_i}$
переводит код каждой точки в $\bar\alpha=(1-\alpha_1,\dots,1-\alpha_n)$.
Поэтому по (L)
\[
f^*(\alpha)=\ov{f(\bar\alpha)} .
\]
Наборы $\{0,1\}^n$ разбиваются на $2^{n-1}$ пар противоположных
$\{\alpha,\bar\alpha\}$ (набор не совпадает со своим дополнением).
\end{remarkbox}

%----------------------------------------------------------------------
\begin{solution}[1 (\S\,3)]
\emph{Ответ: формально да "--- это $\ov A$ и $\ov B$, но обе они
по существу одноместны; самодвойственной операции, существенно
зависящей от обоих аргументов, нет.}

По (L) самодвойственность $F\equiv F^*$ равносильна
\begin{equation*}\tag{C}
F(\ov{A_1},\dots,\ov{A_n})\equiv\ov{F(A_1,\dots,A_n)},
\qquad\text{т.\,е.}\qquad f(\bar\alpha)=\ov{f(\alpha)}\ \ \forall\alpha .
\end{equation*}
При $n=2$ пары противоположных наборов: $\{00,11\}$ и $\{01,10\}$.
Свободно выбираются $f(00)$ и $f(01)$, остальное определено:
\begin{center}\small
\begin{tabular}{ccccl}
\toprule
$f(00)$ & $f(01)$ & $f(10)$ & $f(11)$ & формула\\ \midrule
0&0&1&1& $A$\\ 0&1&0&1& $B$\\ 1&0&1&0& $\ov B$\\ 1&1&0&0& $\ov A$\\ \bottomrule
\end{tabular}
\end{center}
Всего четыре самодвойственные двухместные формулы, и ни одна не зависит
от обоих аргументов. Это и есть <<простой перебор>>, упомянутый в упр.~4.
\end{solution}

%----------------------------------------------------------------------
\begin{solution}[2 (\S\,3)]
\textbf{Разминка.}
\begin{center}\small
\begin{tabular}{lll}
\toprule
$F$ & $F^*$ & самодвойственна?\\ \midrule
$A_k$, $\ov{A_k}$ & те же & да\\
$\emptyset$ & $I$ & нет\\
$A\cap B$ & $A\cup B$ & нет (при $A=\emptyset$, $B=I$)\\
$(A\cap B)\cup(\ov A\cap I)$ & $(A\cup B)\cap(\ov A\cup\emptyset)$ & нет: $\ov A\cup B$ и $\ov A\cap B$\\
$(A\cap B)\cup(A\cap C)\cup(B\cap C)$ & $(A\cup B)\cap(A\cup C)\cap(B\cup C)$ & да "--- (1.6)\\
\bottomrule
\end{tabular}
\end{center}
(Проверка (L) на четвертой строке: при подстановке $A\to\ov A$ черта над
$A$ в формуле дает $\ov{\ov A}=A$; дополняются \emph{буквы}, а не
литералы: $F(\ov A,\ov B)=(\ov A\cap\ov B)\cup A=A\cup\ov B$,
$\ov{F(\ov A,\ov B)}=\ov A\cap B=F^*$.)

\textbf{Критерий.} По (C) и (ii) $F$ самодвойственна $\iff$
$f(\bar\alpha)=\ov{f(\alpha)}$ при всех $\alpha$. Наглядно: если
одновременно поменять ролями <<внутри>> и <<снаружи>> у \emph{всех}
исходных множеств, то результат тоже меняет <<внутри>> и <<снаружи>>.
Например, для (1.6) <<точка входит хотя бы в два из трех множеств>> после
дополнения всех трех превращается в <<входит не более чем в одно>>, т.\,е.
в дополнение исходного множества "--- проверка в одну строку.

На языке конституент: дополнение всех аргументов переставляет
$K_\alpha\leftrightarrow K_{\bar\alpha}$, и самодвойственность означает,
что формула из каждой пары $\{K_\alpha,K_{\bar\alpha}\}$ берет \emph{ровно
одну}.

\begin{theorem*}[описание]
$F(A_1,\dots,A_n)$ самодвойственна тогда и только тогда, когда
\[
F\equiv\bigcup_{\alpha\in S}A_1^{\alpha_1}\cap\dots\cap A_n^{\alpha_n},
\]
где $S\subset\{0,1\}^n$ содержит ровно один элемент из каждой пары
$\{\alpha,\bar\alpha\}$. Разным $S$ отвечают разные формулы.
\end{theorem*}
\begin{proof}
<<Только тогда>>: $S=\{\alpha:f(\alpha)=1\}$ и (C). <<Тогда>>:
$f(\alpha)=1\iff\alpha\in S\iff\bar\alpha\notin S\iff f(\bar\alpha)=0$.
Различность "--- по (i).
\end{proof}
Следствия: самодвойственная формула занимает ровно половину диаграммы
Эйлера, $\lambda(F)=2^{n-1}$; их число $\mu(n)=2^{2^{n-1}}$.

\begin{theorem*}[порождение]
$F$ самодвойственна тогда и только тогда, когда
\[
\boxed{\;F\equiv\bigl(A_1\cap G(A_2,\dots,A_n)\bigr)\cup\bigl(\ov{A_1}\cap G^*(A_2,\dots,A_n)\bigr)\;}
\]
для некоторой \emph{произвольной} формулы $G$ от $n-1$ аргументов;
соответствие $G\mapsto F$ взаимно однозначно.
\end{theorem*}
\begin{proof}
Для любой формулы верно разложение Шеннона
$F\equiv(A_1\cap F_1)\cup(\ov{A_1}\cap F_0)$, $F_1=F(I,A_2,\dots)$,
$F_0=F(\emptyset,A_2,\dots)$. Подставим его в (C). Слева
$(\ov{A_1}\cap F_1(\ov{A'}))\cup(A_1\cap F_0(\ov{A'}))$, где
$A'=(A_2,\dots,A_n)$; справа
$(A_1\cap\ov{F_1})\cup(\ov{A_1}\cap\ov{F_0})$. Множество вида
$(A_1\cap X)\cup(\ov{A_1}\cap Y)$ однозначно определяет $X$ и $Y$ (при
$A_1=I$ и $A_1=\emptyset$), поэтому
$F_1(\ov{A'})=\ov{F_0(A')}$, т.\,е. $F_0=\ov{F_1(\ov{A'})}\overset{(L)}{=}F_1^*$.
Обратно, для $F$ указанного вида по (L) $G(\ov{A'})=\ov{G^*(A')}$,
$G^*(\ov{A'})=\ov{G(A')}$, и
$F(\ov A)=(\ov{A_1}\cap\ov{G^*})\cup(A_1\cap\ov G)=\ov{F(A)}$.
Наконец, $G=F(I,A_2,\dots,A_n)$ восстанавливается по $F$.
\end{proof}
Это второй вывод $\mu(n)=2^{2^{n-1}}$: самодвойственных формул от $n$
аргументов столько же, сколько вообще формул от $n-1$ аргумента.

\textbf{Проверка при $n=3$.} Перебирая 16 формул $G(B,C)$, получаем
$F=(A\cap G)\cup(\ov A\cap G^*)$:
\begin{center}\small
\begin{tabular}{lll}
\toprule
$G$ & $G^*$ & $F$\\ \midrule
$\emptyset$ & $I$ & $\ov A$\\
$I$ & $\emptyset$ & $A$\\
$B$, $\ov B$, $C$, $\ov C$ & те же & $B$, $\ov B$, $C$, $\ov C$\\
$B^b\cap C^c$ & $B^b\cup C^c$ & $\maj(\ov A,B^b,C^c)$\\
$B^b\cup C^c$ & $B^b\cap C^c$ & $\maj(A,B^b,C^c)$\\
$B\sd C$ & $\ov{B\sd C}$ & $\ov{A\sd B\sd C}$\\
$\ov{B\sd C}$ & $B\sd C$ & $A\sd B\sd C$\\ \bottomrule
\end{tabular}
\end{center}
Здесь $b,c\in\{0,1\}$, $\maj(X,Y,Z)=(X\cap Y)\cup(X\cap Z)\cup(Y\cap Z)$.
Итого 16 формул: 6 вырожденных ($A,\ov A,B,\ov B,C,\ov C$), 8
<<медиан>> $\maj(A^a,B^b,C^c)$ (среди них (1.6)) и 2 <<четностных>>
$A\sd B\sd C$, $\ov{A\sd B\sd C}$.

\textbf{Итог.} Самодвойственные формулы "--- в точности те, что
дополняют результат при дополнении всех аргументов:
$F(\ov{A_1},\dots,\ov{A_n})\equiv\ov{F(A_1,\dots,A_n)}$. Их можно выписать
(1) через конституенты, выбирая по одной из каждой пары
$\{K_\alpha,K_{\bar\alpha}\}$, или (2) как
$(A_1\cap G)\cup(\ov{A_1}\cap G^*)$ с произвольной $G$.

\textbf{Замечания.} (а) \emph{Замкнутость.} Если $F$ и $G_1,\dots,G_n$
самодвойственны, то и $F(G_1,\dots,G_n)$ самодвойственна; вместе
с проекциями $A_k$ это замкнутый класс (класс $S$ в теории Поста) "---
ср. \S\,4. (б) \emph{Связь с $R$-функциями.} При переходе к функциям
($A_i\leftrightarrow\{\omega_i\ge0\}$) дополнение множества "--- это смена
знака $\omega\mapsto-\omega$, и (C) превращается в
$f(-\omega_1,\dots,-\omega_n)=-f(\omega_1,\dots,\omega_n)$:
самодвойственным операциям отвечают \emph{нечетные} $R$-функции.
Например, $\tfrac12(x+y\mp|x-y|)$ ($\min$ и $\max$) двойственны, но не
самодвойственны, а медиана $\operatorname{med}(x,y,z)$ нечетна, как
и должно быть для (1.6).
\end{solution}

%----------------------------------------------------------------------
\begin{solution}[3 (\S\,3)]
\emph{Ответ: $F\cap F^*\equiv\emptyset$ $\iff$
$F(\ov{A_1},\dots,\ov{A_n})\equiv F(A_1,\dots,A_n)$ (<<четные>> формулы);
для них всегда $F^*\equiv\ov F$ и потому $F\cup F^*\equiv I$; более того,
условие $F\cup F^*\equiv I$ выделяет тот же класс.}

\emph{Перевод на язык функций.} По (i) $F\cap F^*\equiv\emptyset$
$\iff$ $f(\alpha)\wedge f^*(\alpha)=0$ при всех $\alpha$, а по (ii)
это $f(\alpha)\wedge\ov{f(\bar\alpha)}=0$:
\begin{equation*}\tag{1}
f(\alpha)=1\ \Rightarrow\ f(\bar\alpha)=1\qquad\forall\alpha .
\end{equation*}
Применяя (1) к $\bar\alpha$ вместо $\alpha$ (а $\bar{\bar\alpha}=\alpha$),
получаем обратную импликацию, т.\,е.
\begin{equation*}\tag{2}
f(\alpha)=f(\bar\alpha)\qquad\forall\alpha,
\qquad\text{или}\qquad F(\ov{A_1},\dots,\ov{A_n})\equiv F(A_1,\dots,A_n).
\end{equation*}
Обратно, (2) $\Rightarrow$ (1) очевидно. То же в множествах: $F\cap F^*=\emptyset$
значит $F^*\subseteq\ov F$, т.\,е. (по (L)) $\ov{F(\ov A)}\subseteq\ov{F(A)}$,
$F(A)\subseteq F(\ov A)$ при любых $A$; подставив $\ov A$ вместо $A$,
получаем обратное включение.

\begin{theorem*}
$F\cap F^*\equiv\emptyset$ тогда и только тогда, когда $F$ "---
объединение \emph{целых пар} конституент:
$F\equiv\bigcup_{\alpha\in T}(K_\alpha\cup K_{\bar\alpha})$, где $T$ "---
произвольное множество представителей пар. Таких формул $2^{2^{n-1}}$.
\end{theorem*}

\emph{Замкнутая форма.} При $n\ge2$ четные формулы "--- это в точности
\[
\boxed{\;F\equiv H\bigl(A_1\sd A_2,\ A_1\sd A_3,\ \dots,\ A_1\sd A_n\bigr)\;}
\]
с \emph{произвольной} формулой $H$ от $n-1$ аргумента. Действительно,
$\ov{A_1}\sd\ov{A_i}\equiv A_1\sd A_i$, так что такая $F$ четна. Обратно,
для четной $F$ положим $H(B_2,\dots,B_n)=F(\emptyset,B_2,\dots,B_n)$. Если
код точки $x$ есть $\alpha$ и $\alpha_1=0$, то $x\in A_1\sd A_i\iff x\in A_i$,
и правая часть содержит $x$ тогда же, когда $f(\alpha)=1$; если
$\alpha_1=1$, то $x\in A_1\sd A_i\iff x\notin A_i$, код равен
$(0,\bar\alpha_2,\dots,\bar\alpha_n)=\bar\alpha$, и по (2)
$f(\bar\alpha)=f(\alpha)$. Неформально: четная формула <<не видит>> самих
множеств "--- только то, \emph{совпадают или различаются} состояния точки
по отношению к разным множествам.

\emph{Второй вопрос.} Для четной $F$ по (ii) $f^*(\alpha)=\ov{f(\bar\alpha)}=\ov{f(\alpha)}$,
т.\,е. $F^*\equiv\ov F$; значит $F\cup F^*\equiv I$ и $F\cap F^*\equiv\emptyset$
одновременно. Обратно, $F\cup F^*\equiv I$ означает
$f(\alpha)\vee\ov{f(\bar\alpha)}=1$, т.\,е. $f(\bar\alpha)=1\Rightarrow f(\alpha)=1$,
и тем же приемом получаем (2). Итак, для любой $F$ равносильны:
\[
F\cap F^*\equiv\emptyset\iff F\cup F^*\equiv I\iff F^*\equiv\ov F\iff
F(\ov{A_1},\dots,\ov{A_n})\equiv F(A_1,\dots,A_n).
\]
Для произвольной формулы ни одно из первых двух условий не влечет
другое <<по общим соображениям>>: например, для $F=A\cap B$ имеем
$F\cap F^*=A\cap B\ne\emptyset$ и $F\cup F^*=A\cup B\ne I$.

\textbf{Примеры.} $n=1$: четные формулы $\emptyset$ и $I$ ($A$ и $\ov A$
не подходят: $A\cap A^*=A$). $n=2$: $\emptyset$, $A\sd B$,
$\ov{A\sd B}=(A\cap B)\cup(\ov A\cap\ov B)$, $I$ (т.\,е. $H=\emptyset$,
$B_2$, $\ov{B_2}$, $I$); например,
$(A\sd B)^*=(A\cup\ov B)\cap(\ov A\cup B)=\ov{A\sd B}$. $n=3$: 16 формул
$H(A\sd B,A\sd C)$; при $H=B_2\cap B_3$ получаем
$F=(A\sd B)\cap(A\sd C)=K_{100}\cup K_{011}$, и
$F^*=(A\cup\ov B\cup\ov C)\cap(\ov A\cup B\cup C)=\ov F$.

\textbf{Связь с упр.~2.} Самодвойственные формулы ($f(\bar\alpha)=\ov{f(\alpha)}$,
<<нечетные>>) и четные ($f(\bar\alpha)=f(\alpha)$) "--- два крайних
случая: у первых $F\cap F^*=F$, у вторых $F\cap F^*=\emptyset$. Классы
равномощны, и $F$ самодвойственна $\iff$ $F\sd A_1$ четна.
\end{solution}

%----------------------------------------------------------------------
\begin{solution}[4 (\S\,3)]
\emph{Ответ: строго $n$-местные самодвойственные формулы существуют при
$n=1$ и при всех $n\ge3$; единственное исключение "--- $n=2$.}

\emph{Переформулировка.} Будем говорить, что $F$ \emph{существенно
зависит} от $A_j$, если найдутся $A_1,\dots,A_n$ и $A_j'$ с
$F(\dots,A_j,\dots)\ne F(\dots,A_j',\dots)$. Строгая $n$-местность "---
это существенная зависимость от каждого аргумента: если $F$ не зависит
от $A_j$, то $F\equiv F(A_1,\dots,\emptyset,\dots,A_n)$ "--- формула от
$n-1$ аргумента; и наоборот.

\emph{$n=2$} "--- см. решение упр.~1: все четыре самодвойственные формулы
одноместны. (Самодвойственных формул от двух аргументов $2^{2^1}=4$, ровно
столько же, сколько одноместных самодвойственных $A,\ov A,B,\ov B$ "---
<<места не остается>>.)

\emph{$n=3$}: формула (1.6) $M(A,B,C)$ самодвойственна; при $B=I$,
$C=\emptyset$ она равна $A$, и по симметрии существенно зависит от
каждого аргумента.

\emph{Общий случай $n\ge3$.} По теореме о порождении (упр.~2) при
$G=A_2\cap\dots\cap A_n$, $G^*=A_2\cup\dots\cup A_n$ получаем
самодвойственную формулу
\begin{equation*}\tag{4}
F_n=(A_1\cap A_2\cap\dots\cap A_n)\cup\bigl(\ov{A_1}\cap(A_2\cup\dots\cup A_n)\bigr)
\end{equation*}
(<<точка лежит в $A_1$ и во всех остальных множествах, либо не лежит в
$A_1$, но лежит хотя бы в одном из остальных>>; при $n=3$ это
$M(\ov A,B,C)$). Прямая проверка: при $P=A_2\cap\dots\cap A_n\subseteq U=A_2\cup\dots\cup A_n$
\[
F_n^*=(A_1\cup U)\cap(\ov{A_1}\cup P)=(A_1\cap P)\cup(\ov{A_1}\cap U)\cup(U\cap P)
=(A_1\cap P)\cup(\ov{A_1}\cap U)=F_n,
\]
так как $U\cap P=P=(A_1\cap P)\cup(\ov{A_1}\cap P)$ поглощается.
Существенная зависимость: при $A_2=I$, $A_3=\dots=A_n=\emptyset$ (нужен
аргумент $A_3$, т.\,е. $n\ge3$) $F_n=\ov{A_1}$; при $A_1=I$ и всех
$A_i=I$, кроме $A_j$ ($j\ge2$), $F_n=A_j$. При $n=2$ та же формула
вырождается в $(A_1\cap A_2)\cup(\ov{A_1}\cap A_2)=A_2$.

\emph{$n=1$}: $A$ (и $\ov A$) самодвойственна и не равна тождественно ни
$\emptyset$, ни $I$.

\emph{Оценка подсчетом.} Нестрогая самодвойственная формула равна самодвойственной формуле от
$n-1$ аргумента, поэтому их не больше $n\,\mu(n-1)$, тогда как
$\mu(n)=\mu(n-1)^2$. При $n=3$: $16>3\cdot4=12$, при $n=4$: $256>64$;
при $n=2$ оценка дает $2\mu(1)=4=\mu(2)$ "--- запаса нет, и перебор
показывает, что строго двухместных действительно нет.
\end{solution}

%----------------------------------------------------------------------
\begin{solution}[5 (\S\,3)]
\emph{Что считается.} Различных \emph{записей} бесконечно много
($A_1$, $A_1\cap A_1$, \dots), но формул с точностью до тождества
$2^{2^n}$ "--- по числу таблиц $f$ (см. (i)). Самодвойственные среди них
и составляют $\mu(n)$. При этом <<$n$-местная>> формула не обязана
существенно зависеть от всех аргументов.

\emph{Монотонность без вычислений.} Всякая самодвойственная формула
$F(A_1,\dots,A_n)$ является самодвойственной и как формула от
$A_1,\dots,A_n,A_{n+1}$ (дополнение лишнего аргумента ни на что не
влияет), а различные формулы остаются различными. Значит,
$\mu(n)\le\mu(n+1)$.

\emph{Точное значение.} По критерию (C) таблица задается выбором значения
на одном наборе каждой из $2^{n-1}$ пар $\{\alpha,\bar\alpha\}$, выборы
независимы:
\[
\mu(n)=2^{2^{n-1}}:\qquad
\begin{array}{c|cccccc}
n&1&2&3&4&5&6\\\hline
\mu(n)&2&4&16&256&65\,536&4\,294\,967\,296
\end{array}
\]
Отсюда $\mu(n+1)=\mu(n)^2>\mu(n)$ "--- функция даже строго возрастает.
Для сравнения: всех $n$-местных формул $2^{2^n}=\mu(n)^2=\mu(n+1)$,
т.\,е. самодвойственных формул от $n$ аргументов ровно столько, сколько
вообще формул от $n-1$ аргумента.
\end{solution}

%----------------------------------------------------------------------
\begin{solution}[6 (\S\,3)]
\emph{Ответ: $\tau(2)=0$, и ни при каком другом $n\ge1$ $\tau(n)$ не
обращается в нуль: $\tau(1)=2$, $\tau(3)=10$, $\tau(4)=208$,
$\tau(5)=64\,386$, \dots}

$\tau(2)=0$ "--- решение упр.~1; $\tau(1)\ge1$ ($A$) и $\tau(n)\ge1$ при
$n\ge3$ (формула (4) из решения упр.~4).

\emph{Точная формула.} У каждой формулы есть однозначно определенное
множество $S\subseteq\{1,\dots,n\}$ существенных аргументов; подставив
$\emptyset$ вместо несущественных, получаем строго $|S|$-местную формулу
$\Phi$, которая самодвойственна вместе с $F$. Обратно, пара <<$S$ +
строго $|S|$-местная самодвойственная $\Phi$>> дает самодвойственную
$n$-местную формулу, и разные пары дают разные формулы. Поэтому
\[
\mu(n)=\sum_{k=0}^{n}\binom nk\tau(k),\qquad \tau(0)=0
\]
(формулы $\emptyset$, $I$ от нуля аргументов не самодвойственны). Отсюда
$\tau(n)=\mu(n)-\sum_{k<n}\binom nk\tau(k)$ (или, обращением,
$\tau(n)=\sum_k(-1)^{n-k}\binom nk\mu(k)$):
\[
\begin{aligned}
\tau(1)&=2,\qquad \tau(2)=4-2\cdot2=0,\qquad \tau(3)=16-3\cdot0-3\cdot2=10,\\
\tau(4)&=256-4\cdot10-6\cdot0-4\cdot2=208,\\
\tau(5)&=65\,536-5\cdot208-10\cdot10-10\cdot0-5\cdot2=64\,386 .
\end{aligned}
\]
(Для $n\le4$ значения проверены прямым перебором.) При $n=3$ десять
формул "--- восемь <<медиан>> $\maj(A^{\pm},B^{\pm},C^{\pm})$ и две
<<четности>>.

\emph{Почему именно $n=2$.} Вырожденных самодвойственных формул не
больше $n\,\mu(n-1)$, а $\mu(n)/(n\,\mu(n-1))=\mu(n-1)/n=2^{2^{n-2}}/n$.
При $n\ge3$ это больше единицы и быстро растет, а при $n=2$ равно
$\mu(1)/2=1$: двойка "--- единственное $n$, при котором возведение
в квадрат еще не обгоняет умножение на $n$ ($2^2=2\cdot2$). Доля строго
$n$-местных среди самодвойственных быстро стремится к единице:
$\tau(4)/\mu(4)=208/256$, $\tau(5)/\mu(5)\approx0{,}98$.
\end{solution}

%----------------------------------------------------------------------
\begin{solution}[7 (\S\,3)]
\emph{Ответ: $\lambda(F)=\lambda(F^*)$ $\iff$ $\lambda(F)=2^{n-1}$, т.\,е.
$F$ составлена ровно из половины конституентных областей (равносильно:
$\lambda(F)=\lambda(\ov F)$). Таких формул $\binom{2^n}{2^{n-1}}$.}

$\lambda(F)$ "--- число конституент, из которых состоит $F$ на полной
диаграмме Эйлера (все $2^n$ областей присутствуют), т.\,е. число единиц
в таблице $f$. Например, при $n=2$: $\lambda(A\cap B)=1$,
$\lambda(A\cup B)=3$, $\lambda(A)=2$, $\lambda(\emptyset)=0$,
$\lambda(I)=4$.

По (ii) $f^*(\alpha)=1\iff f(\bar\alpha)=0$: области $F^*$ "--- это
наборы, противоположные наборам, \emph{не} вошедшим в $F$. Отображение
$\alpha\mapsto\bar\alpha$ "--- перестановка, поэтому
\[
\boxed{\;\lambda(F)+\lambda(F^*)=2^n\;}
\]
(геометрически: $F^*$ "--- центральное отражение дополнения $\ov F$ на кубе
кодов). Значит, $\lambda(F)=\lambda(F^*)\iff\lambda(F)=2^{n-1}$; назовем
такие формулы \emph{уравновешенными}. Любое множество $S$ из $2^{n-1}$
наборов дает уравновешенную формулу, разные $S$ "--- разные формулы:
\[
\binom{2^n}{2^{n-1}}:\quad n=1:\ 2,\quad n=2:\ 6,\quad n=3:\ 70,\quad n=4:\ 12\,870,\ \dots
\]

\textbf{Примеры.} $n=2$: $A$, $\ov A$, $B$, $\ov B$, $A\sd B=K_{01}\cup K_{10}$,
$\ov{A\sd B}=K_{00}\cup K_{11}$. $n=3$: медиана (1.6) "--- 4 области из 8;
$A\sd B\sd C$ "--- 4 области; $A\cap(B\cup C)$ "--- 3 области "--- \emph{не}
уравновешена, ее двойственная $A\cup(B\cap C)$ занимает 5.

\textbf{Связь с упр.~2 и 3.} Самодвойственные формулы берут ровно одну
область из каждой пары $\{K_\alpha,K_{\bar\alpha}\}$, поэтому \emph{все}
уравновешены. Четные берут пары целиком, $m$ пар дают $\lambda=2m$
($0\le m\le2^{n-1}$; крайние случаи $\emptyset$ и $I$); уравновешены
лишь те, что составлены ровно из половины пар, $m=2^{n-2}$, их
$\binom{2^{n-1}}{2^{n-2}}$ (при $n=2$ "--- $A\sd B$ и $\ov{A\sd B}$,
но не $\emptyset$, $I$; при $n=3$ "--- 6 из 16). Уравновешенная формула не
обязана быть ни той, ни другой: при $n=3$, $S=\{000,111,001,010\}$,
\[
F=K_{000}\cup K_{111}\cup K_{001}\cup K_{010}=\ov{A\sd(B\cap C)}
\]
берет пару $\{000,111\}$ целиком (не самодвойственна), а пару
$\{001,110\}$ "--- наполовину (не четна), но $\lambda(F)=4=\lambda(F^*)$.
При $n=3$ из 70 уравновешенных формул 16 самодвойственных, 6 четных
и 48 <<прочих>>.
\end{solution}


\section{Композиция и суперпозиция. Замкнутые множества функций}

Для задания функций в виде формул чаще всего используют так называемые
элементарные функции, которые строятся по правилам образования сложных
функций с помощью системы
\begin{equation}\tag{1.9}
\begin{gathered}
H_{\text{э}}=\Bigl\{x\pm y;\ xy;\ \frac xy;\ x^n;\ a^x\ (a>0);\ \log_ax\ (a>0);\\
\sin x;\ \cos x;\ \tg x;\ \ctg x;\ \arcsin x;\ \arccos x;\ \arctg x;\\
\operatorname{arcctg}x;\ a,\ \forall a\in\R\Bigr\},
\end{gathered}
\end{equation}
состоящей из арифметических операций, основных элементарных функций
и констант.

Примерами элементарных функций являются $e^{2\sin^2xy}$;
$\dfrac{x_1x_2+\arctg x_3}{1+x_1x_2x_3}$; $(2(x_1x_2+x_3)+x_4)+2$;
$\sqrt{t^2}$ и~др. Последняя из них $\sqrt{t^2}=|t|$ ($t$ "---
действительная переменная), если ее умело использовать, позволяет считать
элементарными и многие кусочно-элементарные функции (но не все, если речь
идет о функциях более чем одного аргумента!), т.\,е. функции, задаваемые
различными элементарными формулами в различных подобластях пространства
аргументов. Например, можно убедиться, что функция
\begin{equation}\tag{1.10}
h_1(x)=\begin{cases}0&\text{при }|x|>1;\\ 1+x&\text{при }-1\le x\le0;\\ 1-x&\text{при }0<x\le1\end{cases}
\end{equation}
может быть также задана одной формулой:
\[
h_1(x)=\tfrac12\bigl(1-|x|+|1-|x||\bigr).
\]
(График функции $h_1(x)$ показан на рис.~8.)

\begin{figure}[h]
\centering
\begin{tikzpicture}[scale=1.1]
 \draw[->] (-1.7,0)--(1.8,0) node[below]{$x$};
 \draw[->] (0,-0.3)--(0,1.4) node[left]{$h_1$};
 \draw[thick] (-1.6,0)--(-1,0)--(0,1)--(1,0)--(1.6,0);
 \node[below] at (-1,0){\small $-1$}; \node[below] at (1,0){\small 1};
 \node[below right] at (0,0){\small 0}; \node[left] at (0,1){\small 1};
\end{tikzpicture}
\figcap{Рис. 8.}
\end{figure}

Во времена Ньютона и Лейбница понятие функции по существу исчерпывалось
элементарными функциями. Лишь позже, когда было установлено, что некоторые
интегралы от элементарных функций нельзя представить в элементарном виде,
не все уравнения с элементарными коэффициентами можно разрешить
в элементарных функциях и так далее, возникла общая
теоретико-множественная
\bookpage{23}%
точка зрения и элементарные функции стали рассматриваться как примеры
функций вообще.

По мере развития математики число функций, для которых вводилась
стандартная символика, постепенно расширялось, причем происходило это
в рамках как элементарных функций (например, гиперболические функции
$\operatorname{sh}x$, $\operatorname{ch}x$, $\operatorname{th}x$ и~др.),
так и специальных функций (Эйлера, Бесселя, Матье и~т.\,д.). Заметим, что
расширение символики осуществлялось путем своего рода <<естественного
отбора>> и определялось прежде всего практической значимостью
рассматриваемых функций.

Необходимость совершенствования конструктивных средств математики, которые
до недавнего времени почти исчерпывались базисной системой (1.9), особенно
возросла в последние годы, когда широкое распространение получили
быстродействующие вычислительные машины и возникла проблема повышения
эффективности их использования. В частности, появились серьезные основания
для расширения системы (1.9) как в рамках самих элементарных функций, так
и за их пределами.

В дальнейшем наряду с системой (1.9) будут рассматриваться и другие
базисные системы (например, системы булевых функций, $R$-функций и~др.),
поэтому целесообразно подойти к вопросу формирования сложных функций
с помощью базисных систем с общих позиций.

Пусть $\mathfrak X$ "--- некоторое множество, $\mathfrak X^n$ "--- его $n$-я
степень. Отображения вида
$y=f(x)\colon\mathfrak X^n\to\mathfrak X^m$, $x=(x_1,\dots,x_n)$,
$y=(y_1,\dots,y_m)$ назовем \emph{правильными с алфавитом $\mathfrak X$}.
Множество всех правильных с алфавитом $\mathfrak X$ отображений обозначим
$\mathfrak F(\mathfrak X)$. При $m=1$ правильные отображения вида
$f\colon\mathfrak X^n\to\mathfrak X$ называются $n$-местными ($n$-арными)
\emph{алгебраическими операциями} на $\mathfrak X$~[45].

Пусть $y=f(x)\colon\mathfrak X^n\to\mathfrak X^m$;
$u=\varphi(y)\colon\mathfrak X^m\to\mathfrak X^p$. Отображение
$u=\varphi[f(x)]\colon\mathfrak X^n\to\mathfrak X^p$ принято называть
\emph{композицией} $f$ и $\varphi$. Для композиции применяется также
обозначение $\varphi[f(x)]=\varphi\circ f$.

Введем понятие \emph{суперпозиции}, некоторым образом обобщающее понятие
композиции. Пусть $H$ "--- некоторая система (не обязательно конечная)
правильных с алфавитом $\mathfrak X$ отображений
\begin{equation}\tag{1.11}
H=\{y=\varphi_i(x)\colon\mathfrak X^{n_i}\to\mathfrak X^{m_i},\ i\in I_0\},
\end{equation}
где $I_0$ "--- некоторое множество индексов.

Множество $\M(H)$ суперпозиций системы $H$ включает:

а) отображения, составляющие $H$;

б) тождественное отображение $y\equiv x\colon\mathfrak X\to\mathfrak X$;

в) отображения вида
$f=(\varphi,\psi)\colon\mathfrak X^n\to\mathfrak X^{m+p}$, где
$\varphi\colon\mathfrak X^n\to\mathfrak X^m$,
$\psi\colon\mathfrak X^n\to\mathfrak X^p$,
$(\varphi,\psi)=(\varphi_1,\dots,\varphi_m,\psi_1,\dots,\psi_p)$,
а $\varphi,\psi\in\M(H)$;

г) композиции вида $\varphi\circ\psi$, $\varphi,\psi\in\M(H)$.
\bookpage{24}%

Для того чтобы подчеркнуть происхождение суперпозиций от базисной системы
$H$, элементы множества $\M(H)$ будем называть \emph{$H$-реализуемыми
отображениями}, а если $\mathfrak X$ "--- числовое множество, то
\emph{$H$-реализуемыми функциями}.

Понятие $H$-реализуемой функции является формализацией понятия функции,
заданной аналитически (в виде единого аналитического выражения).
Действительно, если для обозначения функций, составляющих базисную систему
(1.11), ввести некоторый набор символов, то для каждой $H$-реализуемой
функции можно будет написать формулу, содержащую лишь указанные символы,
символы независимых переменных и, возможно, открывающие и закрывающие
скобки.

Классическим примером применения операции суперпозиции являются
элементарные функции, которые в соответствии с введенной терминологией
можно назвать $H_{\text{э}}$-реализуемыми функциями (с базисной системой
(1.9)).

Заметим, что при образовании суперпозиций (в силу пункта <<б>> их
определения) допускается произвольное переименование независимых
переменных. Так, например, суперпозициями системы
$H=\{f(x_1,x_2,x_3),\ \varphi(x_1,x_2),\ \psi(x_1,x_2)\}$ являются функции
$f[\varphi(x_1,x_2),x_2,\psi(x_3,x_4)]$;
$f[\varphi(x_1,x_2),x_1,f(x_1,x_1,x_1)]$; $\psi[\varphi(x,y),\varphi(u,t)]$.

Рассмотрим наряду с системой
$H=\{y=\varphi_i(x)\colon\mathfrak X^{n_i}\to\mathfrak X^{m_i},\ i\in I_0\}$
систему
$H_1=\{y=\psi_i(x)\colon\mathfrak X^{n_i}\to\mathfrak X^{m_i},\ i\in I_0\}$,
полученную заменой
\begin{equation}\tag{1.12}
\begin{pmatrix}\varphi_1,&\varphi_2,&\dots,&\varphi_s,&\dots\\ \psi_1,&\psi_2,&\dots,&\psi_s,&\dots\end{pmatrix}
\end{equation}
функций $\varphi_i$ на $\psi_i$, $i\in I_0$ соответственно. (При этом
функции $\psi_i$ имеют то же число аргументов, что и функции $\varphi_i$.)

Пусть $f\in\M(H)$ "--- некоторая $H$-реализуемая функция. Ей соответствует
некоторая формула $f[\varphi_{\alpha_1},\dots,\varphi_{\alpha_m}]$,
$\alpha_i\in I_0$,\fixed{в оригинале
$f[\varphi_{\alpha_1},\dots,\varphi_{\alpha_n}]$ при последующем
$\varphi_{\alpha_1},\dots,\varphi_{\alpha_m}$.} в которую входят какие-то
функции $\varphi_{\alpha_1},\dots,\varphi_{\alpha_m}$ из $H$. Заменим в этой
формуле, согласно (1.12), функции $\varphi_i$ на $\psi_i$,\fixed{в оригинале
<<Заменим в формуле (1.12) функции $\varphi_i$ на $\psi_i$>>; (1.12) "--- это
таблица замены, а заменяются функции в формуле для $f$.} сохранив при этом
наименование независимых переменных. В результате получим новую функцию
$f_1\in\M(H_1)$, описываемую формулой
$f[\psi_{\alpha_1},\dots,\psi_{\alpha_m}]$.

Формулы $f[\varphi_{\alpha_1},\dots,\varphi_{\alpha_m}]$
и $f[\psi_{\alpha_1},\dots,\psi_{\alpha_m}]$ называются \emph{формулами
одинаковой структуры}.

Пусть, например, $H=H_{\text{э}}$ (1.9), $u=(x_1x_2+3)\,e^{x_1x_2}$, и дана
замена:
\[
\begin{pmatrix}x_1x_2;&x_1+x_2;&e^t;&3\\ x_1+x_2;&x_1x_2;&\sin t;&5\end{pmatrix}.
\]
Тогда формула $v=(x_1+x_2)5+\sin(x_1+x_2)$ имеет ту же структуру, что
и формула $u=(x_1x_2+3)\,e^{x_1x_2}$.

С понятием суперпозиции тесно связано понятие \emph{замкнутого}
множества. Множество $\M_0$ называется \emph{суперпозиционно замкнутым},
если для всякой системы $H\subset\M_0$ множество суперпозиций
$\M(H)\subset\M_0$.
\bookpage{25}%

Если все возможные композиции функций из $\M$ принадлежат $\M$, то
множество $\M$ будем называть \emph{композиционно
замкнутым}.\fixed{в оригинале: <<Если все возможные композиции функций из
$\M$ принадлежат $\M$, но имеются суперпозиции функций из $\M$, не
принадлежащие $\M$, то множество $\M$ будем называть композиционно
замкнутым>>. С оговоркой <<но имеются\dots>> следующая фраза
(<<всякое замкнутое множество является композиционно замкнутым>>) была бы
неверна; оговорка описывает множества, композиционно, но не
суперпозиционно замкнутые.}

В дальнейшем, для сокращения речи, суперпозиционно замкнутые множества
будем называть \emph{замкнутыми}. Очевидно, что всякое замкнутое множество
является композиционно замкнутым.

Примерами замкнутых множеств являются:

1) множество $C^k(\R^n)$ ($n=1,2,\dots$) $k$ раз непрерывно
дифференцируемых везде функций (в том числе множества $C^0=C$ непрерывных
и $C^\infty$ бесконечно дифференцируемых функций);

2) множества $C^R$ целых рациональных функций и $C^Q$ "---
дробно-рациональных функций;

3) множество $C^A$ везде аналитических функций (т.\,е. разложимых в ряд
Тейлора в окрестности любой точки);

4) множество $\mathfrak F(\mathfrak X)$ правильных отображений с непустым
алфавитом $\mathfrak X$;

5) пусть $\M=\mathfrak F(\mathfrak X)$ множество правильных отображений
с алфавитом $\mathfrak X$, а $a\in\mathfrak X$. Рассмотрим множество
$T_a\subset\M$ всех таких отображений
$f\colon\mathfrak X^n\to\mathfrak X^m$ ($\forall n,m\in\N$), которые
переводят элементы вида $(a,a,\dots,a)\in\mathfrak X^n$ в элементы вида
$(a,a,\dots,a)\in\mathfrak X^m$. В частности, если $\mathfrak X\subset\R$
и $a=0\in\mathfrak X$, то отображения вида $f\colon\mathfrak X^n\to\mathfrak X$,
обладающие отмеченным свойством, принято называть \emph{алгебраическими
операциями, сохраняющими~0}. Для таких операций всегда $f(0,0,\dots,0)=0$.
Нетрудно доказать (это непосредственно видно из определения суперпозиции),
что множество $T_a$ замкнуто.\fixed{в оригинале <<множество $\M_a$
замкнуто>>; выше это множество обозначено $T_a$.}

\textbf{Теорема 1.} \emph{Пересечение замкнутых (композиционно
замкнутых) множеств есть замкнутое (композиционно замкнутое) множество.
(Пустое множество считается замкнутым по определению.)}

\textsc{Доказательство.} Пусть $\M_1$ и $\M_2$ "--- замкнутые множества,
а $H\subset\M_1\cap\M_2$. Тогда множество $\M(H)$ $H$-реализуемых функций
должно в силу замкнутости $\M_1$ и $\M_2$ содержаться в каждом из них.
Следовательно, $\M(H)\subset\M_1\cap\M_2$, что и требовалось доказать.

В дальнейшем широко используется понятие полной системы отображений, смысл
которого заключается в следующем. Система $H$ называется \emph{полной} по
отношению к множеству $\M_0$, если $\M_0\subset\M(H)$, т.\,е. если всякая
функция из $\M_0$ $H$-реализуема.

{\small
\noindent\textbf{Пример 1.} Система
$H_0=\{x_1+x_2,\ x_1x_2,\ a,\ \forall a\in\R\}$ полна по отношению
к множеству $C^R$ целых рациональных функций.

\noindent\textbf{Пример 2.} Пусть $H$ "--- множество правильных
непрерывных функций двух аргументов с алфавитом $\mathfrak X=[0,1]$, а $C$
"--- множество правильных непрерывных функций произвольного числа
аргументов с тем же алфавитом. Из решения 13-й проблемы Гильберта,
полученного А.\,Н.~Колмогоровым~[51] и В.\,И.~Арнольдом~[4], следует, что
всякая непрерывная функция $n$ аргументов ($n\ge3$) может быть представлена
как суперпозиция непрерывных функций двух аргументов. Следовательно,
множество $H$ представляет собой полную систему относительно
множества~$C$.\par}
\bookpage{26}%

Если для функций, составляющих полную систему, введена некоторая
символика, то всякую функцию из $\M_0$ можно представить в виде единого
аналитического выражения, записанного с помощью этой символики.

Приведем некоторое обобщение понятия полной системы. Пусть
$T\colon\M_0\to2^{\M_0}$, $2^{\M_0}$ "--- множество всех подмножеств
$\M_0$. Элементы множества $Tx$, соответствующего элементу $x\in\M_0$,
будем называть \emph{похожими на $x$ по признаку $T$}. Множество
$\M_1\subset\M_0$ называется \emph{замкнутым по признаку $T$}, если для
всякой суперпозиции $f$ отображений из $\M_1$ в этом множестве есть
отображения, похожие на $f$ (по признаку $T$).

Пусть $H\subset\M_1$ "--- система отображений такая, что для всех
$x\in\M_1$ множество $Tx\cap\M(H)$, где $\M(H)$ "--- множество
$H$-реализуемых отображений, непустое. В этом случае систему $H$ будем
называть \emph{полной по признаку $T$} или \emph{достаточно полной}.
Нетрудно заметить, что если $T$ "--- каноническая инъекция $\M$ в $2^{\M}$
(т.\,е. всякому $x\in\M$ соответствует подмножество $\M$, состоящее из
этого же элемента $x$), то полнота по признаку $T$ есть полнота $H$ по
отношению к $\M$ в обычном смысле.

{\small
\noindent\textbf{Пример 3.} Будем предполагать, что функция
$\varphi\in C^k(\Omega)$ <<похожа>> по признаку $T_\varepsilon$ на функцию
$f\in C^k(\Omega)$, если $\|f-\varphi\|_{C^k(\Omega)}<\varepsilon$. Система
$H_0=\{x_1+x_2;\ x_1x_2;\ a,\ \forall a\in\R\}$ является полной по
признаку $T_\varepsilon$ в $C^k(\Omega)$, каким бы ни было
$\varepsilon>0$. Это следует из известной теоремы Вейерштрасса~[144]
о возможности приближения функций из $C^k(\Omega)$ полиномами (т.\,е.
элементами $\M(H_0)$).

\noindent\textbf{Пример 4.} Пусть $\M$ "--- множество всех функций,
определенных на отрезке $[0,1]$ и имеющих на нем конечное число нулей.
Каждой функции $f\in\M$ поставим в соответствие множество $T(f)\in\M$
функций, нули которых совпадают с нулями $f$. Тогда достаточно полной
является, например, система $H_0=\{x_1+x_2;\ x_1x_2;\ a,\ \forall a\in\R\}$.
Действительно, полином $(x_1-x)\dots(x_s-x)\in\M(H_0)\cap\M$ и имеет
произвольно заданные нули $x_1,\dots,x_s$.\par}

Нетрудно заметить, что система $H_{\text{э}}$ (1.9) избыточная в том
смысле, что некоторые из входящих в нее функций можно опустить, не изменив
множества $\M(H_{\text{э}})$ всех элементарных функций. Такими функциями,
например, являются $\tg x$, $\ctg x$, $\arctg x$ и некоторые другие. (Для
развития теоретической математики введение специальных символов для
функций $\tg x$, $\ctg x$ и других, вообще, не обязательно, однако легко
представить себе те трудности, которые возникли бы в прикладной
математике, если бы соответствующие символы не были введены.)

\textbf{Определение.} \emph{Полная в множестве $\M_0$ система
$H=\{\varphi_i\}$, $i\in I$, $I$ "--- некоторое множество индексов,
называется базисом в $\M_0$, если после исключения из нее любой непустой
подсистемы $H_1$ получаем не полную в множестве $\M_0$ систему (т.\,е.
$\M(H\setminus H_1)\not\supset\M_0$).}

\emph{Замечание.} До сих пор предполагалось, что образование суперпозиций
осуществляется путем конечного числа действий над рассматриваемой системой
$H$. Между тем в математике широко исполь\-%
\bookpage{27}%
зуются и бесконечные суперпозиции, и, без сомнения, изучению таких
суперпозиций принадлежит большое будущее. (Эта точка зрения высказана
в конце прошлого столетия русским математиком Н.\,В.~Бугаевым~[12].)
Классическими примерами таких суперпозиций являются бесконечные ряды
(суперпозиции операции сложения $x_1+x_2$) и бесконечные произведения
(суперпозиции операции умножения $x_1x_2$). Оба случая "--- это частные
случаи следующей более общей схемы. Даны двухместная функция
$\varphi(x_1,x_2)$ и последовательность (числовая или функциональная)
$y_0,y_1,\dots,y_n,\dots$ Назовем частичными суперпозициями формулы
\begin{equation}\tag{1.13}
\begin{aligned}
S_1&=\varphi(y_0,y_1);\\
S_2&=\varphi(y_0,\varphi(y_1,y_2));\\
&\dots\dots\dots\dots\dots\\
S_n&=\varphi(y_0,\varphi(y_1,\varphi(y_2,\varphi(\dots,\varphi(y_{n-1},y_n)\dots)))),\\
&\dots\dots\dots\dots\dots\dots\dots\dots\dots
\end{aligned}
\end{equation}
При $n\to\infty$ приходим к бесконечной суперпозиции со своими признаками
сходимости и методами вычислений. При этом, если
$\varphi(x_1,x_2)\equiv x_1*x_2$ "--- ассоциативная операция, то можно
записать $S_n=y_0*y_1*\dots*y_n$.

Заметим, что ассоциативные операции далеко не редкость даже среди обычных
элементарных функций. Ассоциативные (и коммутативные одновременно),
например, следующие функции:

$\varphi_0(x_1,x_2)=\dfrac{x_1+x_2}{1+c^{-2}x_1x_2}$ (сложение скоростей
в специальной теории относительности);\fixed{в оригинале
$1+c^{-1}x_1x_2$; закон сложения скоростей
$\frac{u+v}{1+uv/c^2}$ содержит $c^{-2}$ (ассоциативность от значения
постоянной не зависит).}

$\varphi_1(x_1,x_2)=x_1+x_2+\alpha x_1x_2$ (при $\alpha=-1$ "--- сложение
вероятностей независимых событий);

$\varphi_2(x_1,x_2)=\sqrt{x_1^2+x_2^2+a^2}$ ($a\in\R$);

$\varphi_3(x_1,x_2)=\dfrac{x_1x_2}{x_1+x_2-\alpha x_1x_2}$ ($\alpha\in\R$;
при $\alpha=0$ "--- сложение сопротивлений в электротехнике).

Впрочем, для построения интересных в теоретическом и прикладном
отношениях бесконечных суперпозиций коммутативность и ассоциативность
вовсе не обязательны. Так, если взять за основу функцию
$\varphi(x_1,x_2)=x_1+\dfrac1{x_2}$, то процесс образования суперпозиций
(1.13) запишется в виде
\begin{equation}\tag{1.14}
S_1=y_0+\frac1{y_1};\qquad
S_2=y_0+\cfrac1{y_1+\cfrac1{y_2}};\qquad
S_3=y_0+\cfrac1{y_1+\cfrac1{y_2+\cfrac1{y_3}}};\qquad\dots
\end{equation}
\bookpage{28}%
т.\,е. приходим к цепным дробям, интерес к которым с каждым годом
возрастает~[159].

Нетрудно заметить, что при замене в формулах (1.13) функции
$\varphi(x_1,x_2)$ другими функциями двух аргументов, будем получать
формулы той же структуры. (Это означает, в частности, что ряды,
бесконечные произведения и цепные дроби (1.14) имеют одинаковую структуру
в том смысле, который определен в данном параграфе.) Нетрудно заметить,
что, используя функции другого числа переменных или чередуя по тем или
иным циклическим законам две или большее число различных формул, можем
получать другие суперпозиционные бесконечные структуры, которые могут,
в частности, иметь и ветвящийся характер. Например, используя функции
$\varphi(x_1,x_2)$ и $\psi(x_1,x_2)$, можем вместо (1.13) рассмотреть
процесс
\begin{equation}\tag{1.15}
\begin{aligned}
S_1&=\varphi(y_0,y_1);\\
S_2&=\varphi(y_0,\psi(y_1,y_2));\\
S_3&=\varphi(y_0,\psi(y_1,\varphi(y_2,y_3))),\\
&\dots\dots\dots\dots\dots\dots\dots
\end{aligned}
\end{equation}
в котором функции входят в формулы для $S_k$, чередуясь.

Если бесконечная структура, опирающаяся на некоторую систему функций
$\varphi_i(x_{1i},x_{2i},\dots,x_{ni})$ ($i=1,2,\dots,k$), определена, то,
осуществляя различные замены вида
\[
\begin{pmatrix}\varphi_1,&\dots,&\varphi_k\\ \psi_1,&\dots,&\psi_k\end{pmatrix}
\]
и выбирая функции $\{\psi_i\}$ тех же аргументов, что и $\varphi_i$ из
некоторых функциональных множеств $\{\Psi_i\}$, получаем различные
бесконечные суперпозиционные процессы одинаковой структуры со своими
критериями сходимости и методами вычислений. При этом возникает ряд
вопросов проблемного характера, касающихся типов сходимости
и аппроксимационных свойств последовательностей $S_n$ вида (1.13), (1.15)
(это относится к случаю, когда $\{y_i\}=\{y_i(x)\}$ "--- функциональные
последовательности), нахождения тех подмножеств множества
$\Psi_1\times\dots\times\Psi_k$, которым соответствуют одинаковые критерии
сходимости и~др. Изучение таких вопросов особенно актуально в связи
с совершенствованием вычислительной техники в направлении создания ЭВМ
с переменной структурой (в частности, так называемых однородных
вычислительных сред).

\subsection*{Упражнения и задачи}
1. Пусть $x*y\colon\R^2\to\R$ "--- некоторая ассоциативная операция,
$h(x)\colon\R\to\R$ "--- строго монотонная функция, $h^{-1}(x)$ "---
функция, обратная $h(x)$. Показать, что двухместная операция
$\varphi(x,y)\equiv h^{-1}[h(x)*h(y)]$ ассоциативна.

2. Построить ассоциативную, но не коммутативную двухместную операцию
$x*y\colon\R^2\to\R$.

3. Пусть $h_t(x)=(1-t)e^x+tx$, $h_t^{-1}(x)$ "--- обратная функция
$h_t(x)$. Показать, что функция $\varphi_t(x,y)=h_t^{-1}[h_t(x)h_t(y)]$
комму\-%
\bookpage{29}%
тативна и ассоциативна при $0\le t\le1$, а $\varphi_0(x,y)\equiv x+y$,
$\varphi_1(x,y)\equiv xy$.

4. Пусть $A$ "--- функциональное множество, $Q(A)$ "--- наименьшее
замкнутое множество, такое, что $A\subset Q(A)$ ($Q(A)$ "--- замкнутая
оболочка $A$). Показать, что $Q(A)=\M(H)$, где $H$ "--- полная по
отношению к $A$ система функций из $A$. Показать, что при
$H\cap\ov A\ne\emptyset$ будет $Q(A)\subset\M(H)$.\fixed{в оригинале <<при $H\cap\ov A\ne\emptyset$ $Q(A)\subset\M(H)$>> (пропущено связующее слово).} Привести примеры.

5. Пусть $A_i$ ($i=1,\dots,m$) "--- функциональные множества,
$H_i\subset A_i$ "--- соответствующие им полные системы. Показать, что
\[
\bigcup_{i=1}^{i=m}A_i\subset\bigcup_{i=1}^{i=m}\M(H_i)\subset\M\Bigl(\bigcup_{i=1}^{i=m}H_i\Bigr).
\]

6. Функция $\varphi(x_1,x_2)$ определяет ассоциативную двухместную
операцию, если $\varphi(x_1,\varphi(x_2,x_3))\equiv\varphi(\varphi(x_1,x_2),x_3)$.
Аналогично функция $\psi(x_1,x_2,x_3)$ определяет ассоциативную
трехместную операцию, если
\[
\psi(\psi(x_1,x_2,x_3),x_4,x_5)\equiv\psi(x_1,\psi(x_2,x_3,x_4),x_5)\equiv\psi(x_1,x_2,\psi(x_3,x_4,x_5)).
\]
Привести пример трехместной ассоциативной операции, не прибегая
к конструкции $\psi(x_1,x_2,x_3)\equiv\varphi(x_1,\varphi(x_2,x_3))$, где
$\varphi(x_1,x_2)$ "--- двухместная ассоциативная операция. Обобщить эти
рассуждения для $n$-местных операций ($n>3$).

7. Построить базисы для множеств целых рациональных функций с целыми,
рациональными и действительными коэффициентами.

% Решения упражнений § 4 (с. 28--29).

\begin{solution}[1 (\S\,4)]
Чтобы выражение $h^{-1}[h(x)*h(y)]$ имело смысл, нужно, чтобы
$h(x)*h(y)$ лежало в области значений $h$; будем считать, что множество
$h(\R)$ замкнуто относительно $*$ (например, $h$ "--- биекция $\R\to\R$).
Строгая монотонность обеспечивает существование $h^{-1}$ на $h(\R)$
и тождество $h(h^{-1}(u))=u$. Тогда
\[
\varphi(\varphi(x,y),z)=h^{-1}\bigl[h\bigl(h^{-1}[h(x)*h(y)]\bigr)*h(z)\bigr]
=h^{-1}\bigl[(h(x)*h(y))*h(z)\bigr],
\]
\[
\varphi(x,\varphi(y,z))=h^{-1}\bigl[h(x)*h\bigl(h^{-1}[h(y)*h(z)]\bigr)\bigr]
=h^{-1}\bigl[h(x)*(h(y)*h(z))\bigr],
\]
и правые части совпадают по ассоциативности $*$. (Это <<перенос
структуры>>: $h$ "--- изоморфизм операции $\varphi$ на операцию $*$; так же
переносится и коммутативность.)

\emph{Замечание.} Все ассоциативные операции $\varphi_0$--$\varphi_3$ из
текста параграфа имеют такой вид:
\begin{itemize}[nosep]
\item $\varphi_0$: $h(x)=\operatorname{arth}(x/c)$, $*=+$ (сложение
  <<быстрот>>; $\operatorname{th}(a+b)=\frac{\operatorname{th}a+\operatorname{th}b}{1+\operatorname{th}a\operatorname{th}b}$);
\item $\varphi_1$ ($\alpha\ne0$): $h(x)=1+\alpha x$, $*=\times$, так как
  $(1+\alpha x)(1+\alpha y)=1+\alpha(x+y+\alpha xy)$;
\item $\varphi_2$ (на $x\ge0$): $h(x)=x^2+a^2$, $*=+$;
\item $\varphi_3$: $h(x)=1/x-\alpha$, $*=+$, так как
  $1/\varphi_3=1/x_1+1/x_2-\alpha$.
\end{itemize}
\end{solution}

\begin{solution}[2 (\S\,4)]
Простейший пример "--- левая проекция $x*y=x$:
$(x*y)*z=x=x*(y*z)$, но $1*2=1\ne2=2*1$.

Менее тривиальный пример: $x*y=x\,|y|$. Ассоциативность:
$(x*y)*z=x|y|\,|z|$, $x*(y*z)=x\,\bigl|y|z|\bigr|=x|y||z|$. Некоммутативность:
$1*(-1)=1$, $(-1)*1=-1$. По упр.~1 из него получается целое семейство
ассоциативных некоммутативных операций $h^{-1}[h(x)\,|h(y)|]$.
\end{solution}

\begin{solution}[3 (\S\,4)]
При $0\le t\le1$ имеем $h_t'(x)=(1-t)e^x+t>0$, так что $h_t$ строго
возрастает. Область значений: при $t=0$ это $(0,\infty)$, при $0<t\le1$
"--- все $\R$ (при $x\to-\infty$ $h_t\sim tx\to-\infty$, при $x\to+\infty$
$h_t\to+\infty$). В обоих случаях она замкнута относительно умножения,
поэтому по упр.~1 (с $*=\times$) операция
$\varphi_t(x,y)=h_t^{-1}[h_t(x)h_t(y)]$ ассоциативна, а коммутативна,
так как коммутативно умножение.

Крайние значения: $h_0=e^x$, $h_0^{-1}=\ln$,
$\varphi_0=\ln(e^xe^y)=x+y$; $h_1(x)=x$, $\varphi_1=xy$. Таким образом,
семейство $\varphi_t$ непрерывно деформирует сложение в умножение,
оставаясь в классе коммутативных ассоциативных операций.

При $0<t<1$ обратная функция выражается через функцию Ламберта $W$
($we^w=s$): из $y=(1-t)e^x+tx$, полагая $x=y/t-w$, получаем
$we^w=\frac{1-t}{t}e^{y/t}$, т.\,е.
\[
h_t^{-1}(y)=\frac yt-W\Bigl(\frac{1-t}{t}\,e^{y/t}\Bigr).
\]
\end{solution}

\begin{solution}[4 (\S\,4)]
\emph{Существование $Q(A)$.} Множество всех правильных отображений
$\mathfrak F(\mathfrak X)$ замкнуто и содержит $A$, а пересечение любого
семейства замкнутых множеств замкнуто (доказательство теоремы~1 дословно
переносится на любое число множеств). Поэтому $Q(A)$ "--- пересечение всех
замкнутых множеств, содержащих $A$.

\emph{$\M(H)$ замкнуто.} Суперпозиции суперпозиций функций из $H$ сами
являются суперпозициями функций из $H$: $\M(\M(H))=\M(H)$ (пункты
а)--г) определения применяются к уже построенным суперпозициям).

\emph{Равенство $Q(A)=\M(H)$ при $H\subset A$, $A\subset\M(H)$.}
(1) $\M(H)$ замкнуто и содержит $A$, значит, $Q(A)\subset\M(H)$.
(2) $H\subset A\subset Q(A)$ и $Q(A)$ замкнуто, значит,
$\M(H)\subset Q(A)$.

\emph{Случай $H\cap\ov A\ne\emptyset$} (система полна по отношению к $A$,
но содержит функции не из $A$). Рассуждение (1) остается в силе:
$Q(A)\subset\M(H)$; рассуждение (2) "--- нет, и включение может быть
строгим.

\emph{Примеры.} $A=C^R$ (полиномы), $H=H_0=\{x_1+x_2,\ x_1x_2,\ a\}$:
$H\subset A$, $Q(C^R)=\M(H_0)=C^R$ (множество $C^R$ само замкнуто).
$A=\{x_1+x_2,\ x_1x_2\}\cup\R$ (константы): $Q(A)=C^R$ "--- замкнутая
оболочка конечного набора операций может быть гораздо шире самого набора.
$A=C^R$, $H=H_0\cup\{\sin x\}$: $H\cap\ov A=\{\sin x\}$, и
$Q(A)=C^R\subsetneq\M(H)\ni\sin x$.
\end{solution}

\begin{solution}[5 (\S\,4)]
Первое включение "--- определение полноты: $A_i\subset\M(H_i)$ для
каждого $i$. Второе следует из монотонности операции $\M$: если
$H'\subset H''$, то всякая суперпозиция функций из $H'$ есть суперпозиция
функций из $H''$, т.\,е. $\M(H')\subset\M(H'')$; применяем это
к $H_i\subset\bigcup_jH_j$.

Включения, вообще, строгие: при $H_1=\{x_1+x_2\}$, $H_2=\{x_1x_2\}$
функция $x_1x_2+x_3$ принадлежит $\M(H_1\cup H_2)$, но не принадлежит
$\M(H_1)\cup\M(H_2)$.
\end{solution}

\begin{solution}[6 (\S\,4)]
\emph{Пример.} $\psi(x_1,x_2,x_3)=x_1-x_2+x_3$ на $\R$. Все три способа
расстановки дают одно и то же:
\begin{multline*}
\psi(\psi(x_1,x_2,x_3),x_4,x_5)=\psi(x_1,\psi(x_2,x_3,x_4),x_5)={}\\
=\psi(x_1,x_2,\psi(x_3,x_4,x_5))=x_1-x_2+x_3-x_4+x_5 .
\end{multline*}
(В любой группе аналогично $\psi=x_1x_2^{-1}x_3$; это <<груда>> Прюфера
и Бэра.)

\emph{Эта операция не имеет вида $\varphi(x_1,\varphi(x_2,x_3))$
с ассоциативной $\varphi$.} Допустим противное:
$\varphi(x,\varphi(y,z))=\varphi(\varphi(x,y),z)=x-y+z$. Пусть $R$ "---
множество значений $\varphi$. Если $u=\varphi(y,z)\in R$, то
$\varphi(x,u)=x+(z-y)$ при всех $x$; значит, $\varphi(x,u)=x+k(u)$ для
$u\in R$ (величина $k(u)=z-y$ определена однозначно). Аналогично
$\varphi(v,z)=z+m(v)$ для $v\in R$. При $u,v\in R$ отсюда
$v+k(u)=\varphi(v,u)=u+m(v)$, т.\,е. $u-k(u)=v-m(v)=c$ "--- постоянная, и
$\varphi(x,u)=x+u-c$ при $u\in R$. Тогда
$x-y+z=\varphi(x,\varphi(y,z))=x+\varphi(y,z)-c$, т.\,е.
$\varphi(y,z)=z-y+c$ при всех $y,z$; следовательно, $R=\R$ и одновременно
$\varphi(x,u)=x+u-c$ и $\varphi(x,u)=u-x+c$ при всех $x,u$, что
невозможно. Противоречие.

\emph{Обобщение на $n$-местные операции.} Операция $\psi$ от $n$
аргументов \emph{ассоциативна}, если при подстановке $\psi$ в любое из $n$
мест получается одна и та же функция $2n-1$ аргументов:
\begin{multline*}
\psi(\psi(x_1,\dots,x_n),x_{n+1},\dots,x_{2n-1})\equiv
\psi(x_1,\psi(x_2,\dots,x_{n+1}),\dots,x_{2n-1})\equiv\dots\\
\dots\equiv\psi(x_1,\dots,x_{n-1},\psi(x_n,\dots,x_{2n-1})).
\end{multline*}
Тривиальные примеры дают <<итерации>> двухместной ассоциативной операции,
$\psi=x_1*x_2*\dots*x_n$. Нетривиальные: пусть $\zeta$ "--- корень
степени $n-1$ из единицы ($\zeta^{n-1}=1$). Тогда
\[
\psi(x_1,\dots,x_n)=x_1+\zeta x_2+\zeta^2x_3+\dots+\zeta^{n-1}x_n
\]
ассоциативна: при подстановке внутреннего $\psi$ на $k$-е место его
$j$-й аргумент получает множитель $\zeta^{k-1}\zeta^{j-1}=\zeta^{(k+j-1)-1}$
"--- ровно тот, что положен месту $k+j-1$, а аргументы, сдвинувшиеся на
$n-1$ мест вправо, сохраняют множитель, так как $\zeta^{n-1}=1$.
В действительных числах можно взять $\zeta=-1$ при нечетном $n$:
$\psi=x_1-x_2+x_3-\dots+x_n$ (при $n=3$ "--- пример выше). Вообще, если
$h$ "--- биекция, то $h^{-1}[\psi(h(x_1),\dots,h(x_n))]$ снова
ассоциативна (как в упр.~1). Систематически $n$-местные ассоциативные
операции изучаются в теории $n$-арных групп (Дёрнте, 1928; Пост, 1940).
\end{solution}

\begin{solution}[7 (\S\,4)]
Функции из базиса должны принадлежать самому множеству $\M_0$, а для
проверки того, что система $H$ "--- базис, достаточно убедиться, что она
полна и что после удаления \emph{любого одного} элемента полнота
теряется. Заметим главное: при суперпозиции полиномов коэффициенты
результата лежат в кольце, порожденном (сложением и умножением)
коэффициентами исходных полиномов.

\textbf{Целые коэффициенты} ($\Z[x_1,x_2,\dots]$):
\[
H=\{x_1+x_2,\ x_1x_2,\ -1\}.
\]
Полнота: $1=(-1)(-1)$, $0=1+(-1)$, $n=1+\dots+1$, $-n=(-1)\,n$; далее
одночлены "--- произведениями, многочлены "--- суммами. Минимальность:
без $-1$ нет ни одной константы (из $x_1+x_2$ и $x_1x_2$ без констант
получаются только многочлены с неотрицательными коэффициентами без
свободного члена); без $x_1x_2$ получаются лишь функции вида
$c+\sum n_ix_i$ степени $\le1$; без $x_1+x_2$ "--- лишь одночлены
$\pm x_1^{k_1}\cdots x_s^{k_s}$ и константы $\pm1$.

\textbf{Рациональные коэффициенты} ($\Q[x_1,x_2,\dots]$):
\[
H=\{x_1+x_2,\ x_1x_2,\ -1\}\cup\{1/p:\ p\ \text{простое}\}.
\]
Полнота: кольцо, порожденное $-1$ и всеми $1/p$, есть $\Q$
($1/q$ для составного $q$ "--- произведение $1/p$). Минимальность: без
$1/p$ все коэффициенты лежат в кольце $\Z[1/q:\ q\ne p]$, где у
знаменателей нет множителя $p$, так что $1/p$ не получить; без $-1$ все
коэффициенты неотрицательны; без $x_1x_2$ или $x_1+x_2$ "--- как выше.
\emph{Конечного} базиса здесь нет: у конечной системы конечное число
коэффициентов, их знаменатели делят некоторое $N$, и все суперпозиции
имеют коэффициенты из $\Z[1/N]\ne\Q$.

\textbf{Действительные коэффициенты} ($C^R$). Система
$H_0=\{x_1+x_2,\ x_1x_2,\ a,\ \forall a\in\R\}$ полна (пример~1), но
базисом в смысле определения не является: отдельные константы избыточны
(например, $2=1+1$). Если, как в обозначении $H_0$ книги, считать семейство
всех констант одним элементом системы, то $H_0$ "--- базис: без сложения,
без умножения или без констант полноты нет (рассуждения те же, что выше).
Если же требовать неизбыточности поэлементно, то по замечанию о
коэффициентах базис вида $\{x_1+x_2,\ x_1x_2\}\cup C$ существует тогда
и только тогда, когда у кольца $\R$ есть неизбыточная система
образующих $C$ (такая, что ни одно $c\in C$ не выражается через остальные
с помощью сложения и умножения). Для $\Z$ и $\Q$ такие системы указаны
выше ($\{-1\}$ и $\{-1\}\cup\{1/p\}$); для $\R$ этот теоретико-кольцевой
вопрос выходит за рамки книги, и мы оставляем его открытым (в предисловии
автор предупреждает, что часть задач носит проблемный характер).
\end{solution}



\appendix
\chapter{Исправленные опечатки и замечания}

\section*{Опечатки оригинала, исправленные в тексте}
\begin{center}\small
\begin{tabular}{@{}r p{0.40\textwidth} p{0.40\textwidth}@{}}
\toprule
с. & в оригинале & исправлено\\ \midrule
3 & <<\dots в которых \texttt{\}}на находит приложение>> & <<\dots она находит\dots>>\\
7 & <<а в последствии>> & <<а впоследствии>>\\
11 & <<(обобщенном )>> & <<(обобщенном)>>\\
13 & <<прообраз множества $A$ при отображении $|$>> & <<\dots при отображении $f$>>\\
15 & Пример: $(A\cap\ov B\cap C)\cap(A\cap C)\equiv\dots$ & $(A\cap\ov B\cap C)\cup(A\cap C)\equiv\dots$ (иначе не проходит вынесение $A\cap C$ по $5^\circ$)\\
16 & в выкладке упр.~2: <<$([A\cap A)\cap(C\cap\ov C)]$>> & <<$[(A\cap A)\cap(C\cap\ov C)]$>>\\
18 & упр.~1 \S\,2: $A\cup[(B\setminus\ov A)\sd C]\equiv(A\cap B)\sd C$ (не тождество) & $A\cap[(B\setminus\ov A)\sd C]\equiv(A\cap B)\sd(A\cap C)$ (восстановлено; см. решение)\\
19 & доказательство (1.6): сомножитель $(B\cup C\cup C)$ на четвертом месте & $(B\cup C\cup B)$\\
21 & нет номера у упражнения <<Будем предполагать\dots>> & <<4.>>\\
24 & $f[\varphi_{\alpha_1},\dots,\varphi_{\alpha_n}]$ & $f[\varphi_{\alpha_1},\dots,\varphi_{\alpha_m}]$\\
24 & <<Заменим в формуле (1.12) функции $\varphi_i$ на $\psi_i$>> & <<Заменим в этой формуле, согласно (1.12), \dots>>\\
25 & определение композиционной замкнутости с оговоркой <<но имеются суперпозиции\dots, не принадлежащие $\M$>> & оговорка снята (иначе неверна фраза <<всякое замкнутое множество является композиционно замкнутым>>)\\
25 & <<множество $\M_a$ замкнуто>> & <<множество $T_a$ замкнуто>>\\
27 & $\varphi_0=\dfrac{x_1+x_2}{1+c^{-1}x_1x_2}$ & $\dfrac{x_1+x_2}{1+c^{-2}x_1x_2}$ (релятивистский закон сложения скоростей)\\
29 & упр.~4: <<при $H\cap\ov A\ne\emptyset$ $Q(A)\subset\M(H)$>> & <<\dots будет $Q(A)\subset\M(H)$>>\\
\bottomrule
\end{tabular}
\end{center}

Мелкие дефекты набора скана (сбитые индексы, разорванные черты
дополнения, точки вместо запятых) исправлены без отдельных отметок.

\section*{Замечания к исходным рабочим заметкам}
Решения упражнений \S\,3 основаны на заметках из каталога
\texttt{notes/claude/} (файлы \texttt{p21-*.md}). При переносе исправлено:
\begin{itemize}
\item В разминочной таблице заметки к упр.~2 формула
  $(A\cap B)\cup(\ov A\cap I)$ отмечена как самодвойственная (<<обе
  $\equiv\ov A\cup B$>>). Это неверно: $F\equiv\ov A\cup B$, а
  $F^*=(A\cup B)\cap\ov A\equiv\ov A\cap B$ (что, кстати, правильно
  вычислено ниже в той же заметке при проверке леммы (L)). В книге эта пара
  приведена лишь как пример \emph{двойственных} формул.
\item В решении упр.~4 добавлена прямая проверка самодвойственности
  формулы $F_n$ и замечание о вырождении при $n=2$; в упр.~6 "--- явная
  ссылка на конструкцию $F_n$; в остальном решения следуют заметкам.
\end{itemize}

\section*{Объем}
Набраны страницы 1--29 оригинала: с. 1--4 (титул, аннотация,
предисловие), с. 5--12 (введение), с. 13--29 (гл.~1, \S\S\,1--4, до
начала \S\,5 <<Машинный способ задания функций. Коммутативные диаграммы>>).
Страница 29 включена, чтобы завершить упражнения \S\,4.

\end{document}
